\chapter{知识的迁移}
\label{chap:ket-transfer-adaptation}

一个已经训练好的花卉分类器，要怎样帮助另一个植物园完成新的工作？这正是迁移学习要回答的问题：能不能让一个在“来源任务”上花了大量数据和算力学到的模型，少走弯路地解决一个数据稀缺的“目标任务”。新园区的需求可能和原任务差得很远：也许要识别来源训练样本里根本没有的物种，也许不再问“这是什么花”而改问“这株有没有病”，甚至要在照片里把每朵花的位置框出来。原模型学到的颜色、纹理、形状这些底层特征很可能仍然通用——识别月季用得上的“花瓣形状”，判断病害时也用得上；但把原模型原封不动跑一次，多半给不出新任务要的答案。真正要做的，是先找出来源经验和目标问题之间到底哪一点相通，再想清楚目标端还缺哪些证据、允许做哪些计算，最后把这份经验变成一种经得起新数据检验的预测能力。这一章就沿“先认清目标、再准备可迁移的东西、最后在目标端完成适应”这条线展开。

第~\ref{chap:ket-update-maintenance}章关心的是“系统随时间改变时，怎么更新又不丢掉旧能力”；本章换一个角度，只问一件事：一个已经存在的来源模型（或来源样本），怎样帮助一个全新的学习目标。这两件事常常同时发生——学新物种和保住旧物种可以是同一次操作的两项责任，但它们的成功标准不同，必须分开验证，不能因为新物种学得好就默认旧物种没忘。沿用前文的损失、风险和参数更新语言，本章分三步：先确定目标任务的预测语义、能拿到哪些样本和访问权限、它和历史任务是什么关系（13.1）；再看在动手碰目标数据之前，来源侧能预先准备好哪些可复用的东西——表示、语义接口、初始化、先验、更新规则（13.2）；最后才落到目标端，区分三种使用经验的方式：构建显式参照、把示范当上下文读取、或用梯度真正调整参数（13.3）。要特别说明，这不是一条“每一步都必须重新训练”的固定流水线：如果现成模型已经够用，完全可以跳过准备直接使用。

这里先澄清三个最容易混在一起的词，因为它们描述的根本不是同一层的条件。目标端只有几张标注图，说的是“目标缺数据”这件事；而准备阶段能用上许多历史任务，说的是“来源很富有”这件事——两者可以同时成立。少样本学习、零样本学习刻画的是前者：目标到底缺了哪一类学习材料（缺标注、还是连新类的图像都没有）。\term{元学习}（meta-learning，也叫“学会学习”）刻画的则是后者：它利用大量历史任务去训练“学习过程本身”，让一个共享模型在接到新任务的少量证据后，能迅速进入可用状态。元学习学的对象很灵活，可以学一个好表示、学一个好的参数起点、也可以学一套读取证据的方式，它并不限定目标端一定要做梯度更新。正因如此，这两个维度是正交的：同一个方案完全可以“在目标端是少样本学习，同时在准备端采用了元学习”。

\section{可迁移的目标任务}
\label{sec:ket-transfer-targets}

新园区的口头需求（“帮我们认认花、看看病”）必须先翻译成一个能打分、能比较的预测问题，否则无从判断迁移有没有用。而且不同任务“错一次”的含义完全不同：识别物种时，错误就是选错了类别；定位花朵时，光选对类别不够，还要评价框出的位置准不准、这个框和那朵花有没有正确配对——这两类错误显然不能共用一个含义都没说清的分数。为统一语言，设来源任务的预测输入空间、真实答案空间和联合分布为 $X_S,Y_S,\mathcal P_S$，目标任务对应为 $X_T,Y_T,\mathcal P_T$。对目标预测器$f$，用$\ell_T$给每个预测与真实答案之间的差异计分，再在未来目标样本上取平均，得到
\begin{equation}
  R_T(f)=\mathbb E_{(\bm{x},y)\sim\mathcal P_T}
  \bigl[\ell_T(f(\bm{x}),y)\bigr].
  \label{eq:ket-transfer-target-risk}
\end{equation}
其中 $(\bm x,y)\in X_T\times Y_T$ 是一张目标图像及其真实答案，$f(\bm x)$ 可以是类别、分数、概率或检测框这样的结构化输出，只要和 $\ell_T$ 的接口对得上即可，并假定损失的期望存在。这里期望刻意对 $\mathcal P_T$ 取，因为我们关心的是模型在“未来会遇到的目标样本”上平均错多少，而不是回头把来源训练集再考一遍（后者再好也说明不了新园区的表现）。还要注意 $R_T$ 的可比性是有前提的：只有把任务、损失、允许访问的观测或标签、以及预算都固定下来，“更小的 $R_T$”才等价于“目标错误更少”；而且这个数字本身并不告诉你好处究竟来自哪一项来源资产。一旦预测输入分布、输出语义或评价标准变了，这套对应关系就得重新核对。

\subsection{预测目标与输出语义的确定}
\label{sec:ket-transfer-output-semantics}

先比较三种“输入仍是花卉照片、但任务实质不同”的目标，看清迁移时到底什么变了。第一种是新增物种：一张图仍只输出一个类别名，变化的是类别集合变大了，需要学新的分类边界。第二种是改判病害：这变的是标签的含义——同一物种的不同植株可以有完全不同的病害状态，而原来“这是月季”的物种标签里根本不包含健康与否的区分。所以病害任务不是把输出编号改个名字就行：把第 3 类改名“白粉病”，既补不上缺失的病害监督，也补不上辨认病斑所需的视觉特征，模型依旧不知道该看叶片上的什么。

第三种定位任务更彻底：它把预测单位从“一整幅图”换成了“图里的一个个对象”。输出必须同时表达有几朵花、每朵在什么位置、各自是什么类别，训练时还得规定模型给出的哪个框对应标注里的哪朵花（否则多框多花时对错无法核算）。如果只是把输出向量加长几个数，却不规定这些坐标代表什么、对象数量不固定时怎么处理、框和对象用什么损失匹配，那就根本没有定义出一个检测任务。同理，分割以像素区域为单位、生成以输出序列为单位，都各自需要匹配的输出结构和监督信号。结论是：来源表示能复用到哪一层，必须按新任务的预测单位逐个检查，不能默认“特征通用”就万事大吉。

即便回到最简单的单图分类，也得先做一件容易被忽略的事：对齐类别含义，而不是只看编号。令 $C_S,C_T$ 分别为来源、目标实际可能出现的语义类别集合，二者交集 $C_{\cap}=C_S\cap C_T$ 就是共享类。编号相同不等于语义相同：来源按“物种”分类、目标按“园艺品种”分类时，一个同名的“月季”在两边覆盖的范围可能差很多，必须先讲清映射关系。表~\ref{tab:ket-transfer-class-relations} 进一步固定一种常用协议来讨论：训练时只有带标签的源样本和无标签的目标样本，目标标签被严格留到方案锁定之后才用于评价。在这个协议下，类别集合之间的包含关系，直接决定了哪些源类该往目标迁、哪些目标图像应该被拒绝识别。要强调的是，“类别集合关系”和“监督权限”是两个独立维度：目标有没有少量标签、能不能访问来源原始训练样本还是只能读到一个源模型，这些都要单独列清楚。

\begin{table*}[htbp]
  \centering
  \small
  \setlength{\tabcolsep}{3pt}
  \begin{tabularx}{\linewidth}{p{19mm}p{25mm}p{29mm}XX}
    \toprule
    设定 & 类别关系 & 代表性训练访问 & 目标端判断 & 额外风险 \\
    \midrule
    完全共享 & $C_T=C_S$ & 带标签源样本；无标签目标样本 & 在共同类别中分类 & 拍摄条件变化仍可破坏原分类器 \\
    部分领域自适应 & $C_T\subsetneq C_S$ & 同左；训练时不揭示目标标签 & 迁移目标实际具有的类别 & 源中特有类别可能被错误地对齐到目标 \\
    开放集领域自适应 & 典型设定为$C_S\subsetneq C_T$ & 同左；目标私有类不提供类别名 & 共享类分类，目标私有类拒识为未知 & 未知不能被强迫匹配到已知类 \\
    通用领域自适应 & 两侧均可有私有类 & 同左；共享范围预先未知 & 共享类分类，其余目标样本预测输出未知 & 共享范围与未知比例都需要估计 \\
    \bottomrule
  \end{tabularx}
  \caption{同一图像分类协议下的类别关系与预测责任。四行均允许带标签源样本和无标签目标样本；集合关系在语义对齐后判断。开放集一行采用$C_S\subsetneq C_T$的常用设定，通用设定不预定包含关系。预测输出“未知”只承认样本不属于共享已知类，不给出各目标私有类的名称。}
  \label{tab:ket-transfer-class-relations}
\end{table*}

先看 $C_T\subsetneq C_S$（目标只是来源的一部分）的情形：这时来源里那些“目标根本没有”的私有类反而成了噪声，硬把它们迁过去会把目标样本错误地吸向这些多余类别。\term{部分领域自适应}（partial domain adaptation）就是在这种包含关系下，想办法压低源私有类对迁移的干扰。代表性的 PADA（部分对抗领域自适应）从模型当前对目标样本的预测里反推每个源类的权重（老被预测到的源类更可能属于共享部分），再用权重同时调节源分类损失和领域对抗损失。所谓领域对抗训练，是训练一个特征提取器让“源特征”和“目标特征”难以被一个域判别器区分，从而对齐两域分布；加权则进一步限定“让哪些源内容参与这场对齐”。不过权重依赖当前预测，存在自我强化的风险：若某个共享类在早期模型里被漏判，它的权重就可能被错误地一路压低\scite{ket-cao2018partial}

反过来，$C_S\subsetneq C_T$（目标多出来源没有的物种）时，麻烦出在目标侧：若强行把所有目标特征都拉向某个源类，就等于逼着一种全新物种去认领一个它不属于的已知名。\term{开放集领域自适应}（open set domain adaptation）因此允许模型输出一个特殊答案——“未知”，即承认这张图不属于任何已知类。Saito 等人让分类器和特征生成器在目标样本上形成对抗：一部分目标样本被推向已知类、与它们对齐，另一部分则被推开、与已知类分离为未知。正因为引入了“未知”出口，评价时就不能只看共享类的准确率，还必须单独检验未知拒识做得好不好\scite{ket-saito2018openset}

更一般的情形是，事先连哪些类共享都不知道，两侧都可能各有私有类。\term{通用领域自适应}（universal domain adaptation）于是干脆不预设标签集之间的包含关系，要求模型仅凭拿得到的证据，自己判断哪些目标样本属于共享内容、其余的标记为未知。相比前两种已知包含关系的设定，它多了一道“估计共享范围”的难题。名字里的“通用”只是说放宽了对类别集合关系的假设，绝不是承诺对任意新任务都能成功\scite{ket-you2019universal}

处理未知类时要分清三件层层递进、但常被混为一谈的事。第一件是\term{未知拒识}（unknown rejection）：它只表达一个很弱的结论——现有的类别目录没法给这张图一个合适答案，仅此而已，它并不认识这是什么。第二件是把一堆被拒识的图像聚成若干群组，这是在进一步挖掘输入空间里的结构，但聚出来的群可能只是“背景都是绿叶”或“都在温室拍的”，未必对应真实物种。第三件才是真正建立一个新类别：你得先确定每个群组在语义上是什么、补入足以把它和别的类区分开的代表性证据，再在新样本上检验分类规则是否成立。一条现实可行的链条通常是“先聚类、再请专家命名和补标、最后训练分类器”；而一个笼统的“未知”输出，连第一步之后的所有事都还没做。

用一个具体例子把边界钉死。设来源类别是 {菊花、月季、鸢尾}，目标类别是 {月季、鸢尾、新种甲}。真正共享的只有月季和鸢尾，所以菊花不该成为新园区“必须从中选一个”的答案；而新种甲来源里压根没有。如果目标只给了无标签照片、既没有新种甲的文字描述也没有标注，那么模型能合理争取到的最好结果只是“把新种甲拒识为未知”，不可能无中生有地为它建一个有语义、能叫出名字的分类器——要获得“认识新种甲”这个新增能力，后续的描述和标注是必不可少的依据。另外别把两件事搞反：在目标任务里排除菊花的干扰，并不等于要求系统把菊花彻底忘掉；旧园区可能仍要识别菊花，那属于保留问题，应按第~\ref{sec:ket-update-retention-validation}节另行验证。

还要把“语义变化”和“分布变化”区分开，它们能同时发生但性质不同。类别不变、只是换了相机或季节，主要改变的是观测（输入）分布，标签关系大体还在；而把“物种标签”换成“病害标签”，哪怕输入分布一模一样，预测任务也已经是另一个了。因此，单纯把源和目标的输入特征对齐，并不能替你建立新的标签联系。更微妙的是，即使输入空间和标签空间都固定，同样一批看得见的图像，在目标域仍可能对应完全相反的标签规律（比如某种斑纹在源域是病、在目标域是品种特征），这时只靠对齐输入根本无法识别目标规则；这类“不可识别”的反例可参考\BookChapterRef{chap:transfer-learning-theory}{卷3《基础、理论与可信性》的“迁移学习的可学习性”章}。新任务也须先提供连接来源与目标的证据，再使用第~\ref{sec:ket-update-weighting}节的加权或式~\eqref{eq:ket-update-crossmodal}的配对约束，并重新检查其覆盖与判别特征保留条件。本书将常被称为“领域适应”的相关问题统一称为领域自适应。

\subsection{目标信息与适应权限的确定}
\label{sec:ket-transfer-evidence}

对新种甲、新种乙，植物园到底愿意提供什么，直接决定模型能从哪里获得“把两者区分开”的依据，而不同材料给的是完全不同的东西。文字描述（“花瓣细长、叶缘有锯齿”）告诉模型这类花应当具备什么特征；带标签图像把“类别名”和“实际看到的像素”连接起来，让特征落地；无标签的当地照片虽然没有答案，却能告诉模型目标域图像长什么样、各形态出现得多频繁、局部结构如何。所以“只给描述”“每类给几张标注图”“再开放一大批无标签照片”是三种证据量不同的条件，哪怕三者在部署时都只是跑一次前向预测，也不能因此把它们当成相同证据条件下的公平比较。

\term{少样本学习}（few-shot learning）描述的就是目标端只有少量学习用标注这一条件。分类里常规定义 $N$ 个类、每类给 $K$ 个样本，记作 $N$-way $K$-shot，支持样本总数即 $NK$。但这个简写其实漏掉了很多影响难度的信息：有多少查询样本、各类是否平衡、支持图像有没有重复、用了什么增强——比较两个方法时这些都得补齐，否则同样叫“5-way 1-shot”难度可能天差地别。若目标只有类别描述、没有该类的训练图像，通常称为\term{零样本学习}（zero-shot learning）。这里“零”这个字必须说清到底缺的是什么，否则极易被误解：没有新类的图像，仍可能有它的名称和属性描述；没有某个语言任务的示范，仍可能有一句任务指令；而且模型很可能在大规模预训练里早就接触过相关概念。因此不同的“零样本”协议，不能只凭名字就拿来比较它们各自需要的观测与标注。

为了把“目标到底给了什么”写得毫无歧义，可以约定一个目标证据包，它由四部分组成：类别描述 $D_T$、带标签的支持集 $S_T$、被许可使用的无标签集 $U_T$、以及已经到达的交互历史 $H_T$。其中任何一项都可以为空，方案也不被要求必须用尽所有许可材料。比较不同方法时，要固定这个证据包以及各项材料“何时可见”，还要把容易藏起来的投入一并记账：来源预训练用了多少样本、专家整理描述花了多少功夫、检索库里的示范、以及额外用来选模型的标签。否则你看到的“目标风险更低”，可能只是因为某方偷偷用了更多样本或更详尽的描述，而不是它的适应操作本身更聪明。

证据数量对齐后，还得约定“怎么使用待预测的查询图像”，这是另一条常被忽略的协议分界线。\term{归纳式预测}（inductive prediction）在适应结束后，一张张独立处理新图像，当前答案不受其他待预测图像影响。\term{直推式预测}（亦称转导式预测，transductive prediction）则允许把一整批无标签查询放在一起利用——比如观察到这批图像自然聚成两团，就据此调整分类边界。直推式多拿到的是“查询图像之间的结构关系”，但它并没有拿到任何隐藏标签，所以这份额外观测带来的好处不能全算到优化器头上。还要考虑可部署性：真实上线时往往是照片一张张陆续到达，那种依赖“整批查询同时在手”的算法未必能照搬到线上。

除了数据，模型参数的读写权限、状态能否保留、反馈什么时候到达，也都是协议的一部分，因为它们改变了“当下到底能执行怎样的学习”。能不能更新骨干网络、能不能跨请求保存记忆、能不能在作答前先向专家发问，每一项许可都对应一种不同强度的方法。至于如何主动发起查询并校验新观测或标注，属第~\ref{part:efficient-knowledge-use}部分的话题；这里只需把查询次数、费用和反馈到达时刻如实记录。一条因果纪律必须守住：事后才拿到的正确标签，可以用来改进以后的预测，但绝不能倒回去用在“当时还没有反馈”的那些预测上，再把更新后的成绩冒充成原协议下的测试成绩。

在这些边界内，比较才有意义：可以固定标注数 $n$、更新步数 $T$ 和计算预算 $B$，比谁的目标风险低；也可以固定一个风险阈值，比谁达到它所需的投入更少。风险当然越低越好，但成本必须用统一单位计量，并且一定要配一个“从零开始的普通目标训练”作对照，否则说不清迁移到底带来了多少净收益。记账时尤其要把那些不显眼的开销纳入：整理描述、检索示范、挑选提示、以及验证调参。只数反向传播次数是个常见陷阱——像零样本、上下文学习这类根本不更新参数的方案，反向传播为零，但它在编码文本、比对图像上的主要适应费用一点也不少。

\subsection{历史任务与目标任务关系的确定}
\label{sec:ket-transfer-task-splits}

为新地区做准备，光有一大堆图像并不够，更关键的是要有一组足以反映“未来会遇到什么样的学习问题”的历史任务。一个任务 $\mathcal T_k$ 完整规定了四件事：预测输入和输出的语义、联合分布 $\mathcal P_k$、损失函数、以及样本和模型的访问权限。把这些任务看作从某个任务分布 $\Pi$ 中抽样得到的，就出现了一个两层抽样结构：先抽出一个“在某地区区分哪几种花”的问题，再从这个问题内部抽出一张张图像。为什么要强调这两层？因为元学习赌的是任务之间的共性，而共性可能可靠也可能虚假：如果历史任务共同依赖的是“花瓣形状”这种稳定规律，学到的表示就能迁到新地区；但如果它们只是碰巧共享了“都在某种背景下拍摄”，那换个地区这种联系立刻失效。所以任何准备方案都必须说清楚，它预期保留和复用的究竟是哪一层、哪一种规律。

在单个任务内部，也要把“用来学的”和“用来考的”分开。用于适应的\term{支持集}（support set）$S_k$ 提供当前任务的输入—标签对，模型据此形成规则；\term{查询集}（query set）$Q_k$ 则检验这条规则能不能用到没见过的样本上。两者当然可以包含相同的类别，但必须按真实的评价单位严格隔离：不能把同一张原图的不同增强版本一个放进支持、一个放进查询（那样答案几乎可以互相照抄）；若评价单位是植株，就连同一植株的所有照片都要整组隔离。再往上，任务之间还要分成元训练、元验证和最终测试三组任务；如果你的主张是“能迁移到没见过的类”，那这三组任务用到的类别本身也必须互不重叠。

那么训练目标应该怎么写？这里有一个关键转向：要判断哪种共享准备更有利于“未来的适应”，就不能只看它在历史图像上直接预测的损失（那只考记忆，不考适应力），而要看“拿到一个新任务的支持集后，它形成的预测器在该任务查询集上表现如何”。把共享准备记为 $\bm\phi$，把“用支持证据形成任务预测器”这一适应过程记为算子 $A_{\bm\phi}(S_k)$，再用非空 $Q_k$ 上的平均损失 $\hat R_{Q_k}$ 给形成的预测器打分。对“抽任务”和“任务内抽样本”这两层都取期望，就得到元学习最典型的准备目标：
\begin{equation}
  \min_{\bm{\phi}}
  \mathbb E_{\mathcal T_k\sim\Pi}
  \mathbb E_{S_k,Q_k\mid\mathcal T_k}
  \bigl[\hat R_{Q_k}(A_{\bm{\phi}}(S_k))\bigr],
  \label{eq:ket-transfer-meta-objective}
\end{equation}
这个式子最小化在预先选定的共享参数范围内进行，期望存在时，它衡量的就是“任务先经支持集适应、再在查询集上的平均错误”，而非适应前的错误——这正是元学习与普通训练的分野。实际中历史任务有限，只能用采样平均去近似这两层期望，也不保证恰好找到全局最小。要特别注意各层样本的身份：元训练任务的查询标签虽然在任务内部扮演“考题”，但它参与了 $\bm\phi$ 的更新，所以对 $\bm\phi$ 而言它仍是训练数据，不是最终的独立测试。模型结构、步长、停止条件这些超参，交给元验证任务去选；而真正留出的目标任务，其查询答案要一直保密到整个方案冻结之后才用于评分。图~\ref{fig:ket-transfer-data-boundaries} 用“历史侧有虚线反馈、目标侧评分反馈不存在”刻画的，正是这两层样本用途的边界。

\begin{figure*}[htbp]
  \centering
  % Editable schematic, designed at 166 mm width; no measured data.
% Solid arrows: evidence use / prediction. Dashed arrows: training feedback.
\begin{tikzpicture}[
  x=1mm,y=1mm,
  font=\figurefont\footnotesize,
  text=BookInk,
  >={Latex[length=2mm,width=1.4mm]},
  evidence/.style={draw=BookBlue,fill=FigureInput!10,line width=0.8pt,
    text width=24mm,minimum height=12mm,align=center,inner sep=2mm},
  model/.style={draw=BookTeal,fill=FigureModel!10,line width=0.8pt,
    text width=26mm,minimum height=12mm,align=center,inner sep=2mm},
  score/.style={draw=BookAmber,fill=FigureLoss!10,line width=0.8pt,
    text width=28mm,minimum height=12mm,align=center,inner sep=2mm},
  flow/.style={draw=BookInk,line width=0.9pt,-{Latex},rounded corners=1mm},
  feedback/.style={draw=BookAmber,line width=0.9pt,dashed,-{Latex},
    rounded corners=1mm}
]
  \node[anchor=west,font=\figurefont\small] at (2,15) {历史任务：允许跨任务训练反馈};
  \node[evidence] (support-one) at (16,0) {$S_1$：月季、鸢尾\\支持图像与标签};
  \node[model] (adapt-one) at (59,0) {用共享准备$\bm\phi$\\形成任务预测器1};
  \node[score] (query-one) at (101,0) {$Q_1$：其他图像\\查询预测与损失};

  \node[evidence] (support-two) at (16,-29) {$S_2$：菊花、百合\\支持图像与标签};
  \node[model] (adapt-two) at (59,-29) {用共享准备$\bm\phi$\\形成任务预测器2};
  \node[score] (query-two) at (101,-29) {$Q_2$：其他图像\\查询预测与损失};
  \node[model,text width=25mm,minimum height=19mm] (prepare) at (150,-14.5)
    {跨任务更新\\共享准备$\bm\phi$};

  \draw[flow] (support-one.east) -- (adapt-one.west);
  \draw[flow] (adapt-one.east) -- (query-one.west);
  \draw[flow] (support-two.east) -- (adapt-two.west);
  \draw[flow] (adapt-two.east) -- (query-two.west);
  \draw[feedback] (query-one.east) -- (123,0) |- (prepare.west);
  \draw[feedback] (query-two.east) -- (123,-29) |- (prepare.west);
  \node[anchor=west,text=BookMuted] at (2,-43)
    {任务1、2均从同一轮准备出发；更新后进入下一轮历史任务};

  \draw[draw=BookRule,line width=0.8pt] (2,-50) -- (165,-50);
  \node[anchor=west,font=\figurefont\small] at (2,-59)
    {留出目标：方案冻结后，只用支持证据适应};

  \node[evidence] (support-target) at (16,-76)
    {$S_T$：新种甲、乙\\允许的图像与标签};
  \node[model] (adapt-target) at (59,-76)
    {目标任务预测器\\由$S_T$形成};
  \node[model,text width=23mm,fill=FigureOutput!8] (predict-target) at (103,-76)
    {$Q_T$预测输入\\目标预测};
  \node[score,text width=24mm] (score-target) at (151,-76)
    {$Q_T$隐藏答案\\最终评分};
  \draw[flow] (support-target.east) -- (adapt-target.west);
  \draw[flow] (adapt-target.east) -- (predict-target.west);
  \draw[flow] (predict-target.east) -- (score-target.west);
  % Frozen transfer is a one-way route outside the historical task lanes.
  \draw[flow] (prepare.south) -- (150,-47) -- (167,-47)
    -- (167,-64) -- (59,-64) -- (adapt-target.north);
  \node[anchor=east,fill=white,inner sep=1mm,text=BookTeal]
    at (148,-64) {冻结准备结果};
  \node[anchor=west,text=BookMuted] at (2,-92)
    {目标评分没有返回适应或模型选择的箭头};
\end{tikzpicture}

  \caption{共享准备如何接收历史反馈而隔离目标答案。实线表示支持证据、冻结准备和预测的使用，虚线表示历史查询损失更新共享准备。历史任务各自从同一轮准备出发；留出目标只用许可的支持证据形成预测器，查询答案只进入最终评分。图中的花卉类别为教学划分，目标还可包含独立留出的领域变化。}
  \label{fig:ket-transfer-data-boundaries}
\end{figure*}

为什么“多几个任务”比“同一任务多几张图”更能说明可迁移性？Baxter 用归纳偏置学习给出了理论刻画：一个候选假设空间规定了每个任务能去拟合哪些函数，任务内的训练样本帮你在这个空间里选定一个函数，而跨多个任务积累的经验则帮你选定“用哪个空间”这件更上层的事。于是两种数据的作用完全不同：给某个任务多添图像，只能改善这一个任务的估计，并没有为“任务之间的关系”增加任何独立证据；只有增加彼此独立的新任务，才能检验你选出的这套限制是否真的适用于未来任务。当然，这套论证依赖任务之间确有联系、也依赖对候选空间族做容量控制，否则共享准备可能只是把有限几个历史任务背了下来\scite{ket-baxter2000-bias}

检验历史任务的覆盖面时，还要把“类别新不新”和“领域变没变”这两个轴分开。同一台设备、只是换一组物种，那只检验了对未见类的泛化；而从温室高清照换成野外红外影像，则连观测方式和领域都变了，是更强的挑战。Meta-Dataset 纳入多种数据来源，并让每个任务的类别数、支持样本数都变化，避免评价被“固定五分类、每类固定几枪”这一种设定锁死；但即便如此，要理直气壮地声称“跨领域迁移”，仍必须把相应的目标领域整个留出不参与训练\scite{ket-triantafillou2019metadataset}。报告成绩时，变化应分别归因到类别新颖性、领域变化、标注数量和类别本身的区分难度，不能把这些变量搅成一句含义不清的“任务更难了”。

当真实历史任务不够多时，一个自然想法是“人造任务”：通过在历史任务之间做插值来扩充准备材料。具体是混合它们的中间特征；如果两边标签语义共享，可以连同标签一起插值，否则就得先把类别配对好、重新规定任务内部的标签。构造时支持集和查询集仍要分别生成，查询答案绝不能漏进支持输入。但要清醒：插值只是定义了一批新的训练问题，它并没有证明这些人造问题真的覆盖了“新地区、新物种”会带来的变化——造出来的任务再多样，也可能只在历史任务张成的空间里打转。所以这类方法的真实收益，最终仍要靠完全独立的目标任务来检验\scite{ket-yao2021interpolation}

还有一个隐蔽的陷阱：查询成绩很高，也不代表模型真的学会了“看支持集做事”。设想“月季”在所有历史任务里都被固定编号为 1，模型完全可能直接死记这个编号映射，根本不去读本次支持集——它照样能答对，但这不是少样本适应，而是作弊式的记忆。要拆穿这类捷径，可以分别做几个对照：把每个任务的标签重新随机一致编码（打破固定编号）、换掉支持集（看答案是否随之改变）、移除输入里泄露任务身份的线索。Yin 等人则用互信息约束去限制模型参数或查询输入中残留的“任务记忆”，逼迫预测真正依赖当下的支持样本；不过这种约束也不能保证彻底消灭所有任务里的记忆旁路\scite{ket-yin2019memorization}。样本隔离的统一规则与任务层的重复分析，分别见第~\ref{sec:ket-evaluation-splits}、\ref{sec:ket-evaluation-repeats}节。

\section{可迁移知识的准备}
\label{sec:ket-transfer-preparation}

目标任务和历史任务的关系厘清后，准备阶段要回答的问题就很具体了：来源模型里到底还缺哪一种“可复用的计算”，值得我们在动手碰目标数据之前先准备好？这个问题的答案取决于目标端拿什么来：若只有少量花卉图像，最需要的是一个能公平比较它们的好表示；若只有类别描述，就得先有一座把“语义”和“视觉”连起来的桥；若允许做几步优化，就得给优化准备一个好起点、一个好先验、甚至一套好的更新规则；若打算让模型直接读示范而不更新参数，则要准备一个“证据到预测”的条件映射。下面就按这五种功能逐个检查准备内容。要强调它们并不是五类互斥、只能选一的模型，实际方案常常几种共存；也不是每种都得做——如果现成的预训练表示已经足够好，就没必要再专门元训练。只有当“未来任务确实共享某种规律、且这类任务会被反复用到、足以摊薄准备成本”时，专门去学“如何适应”才是划算的。

\subsection{可迁移表示的准备}
\label{sec:ket-transfer-representation}

先从最基础的资产——可迁移表示说起。每类只有几张图时，成败几乎完全取决于来源编码器是否已经把“区分类别所需的线索”保留了下来。记编码器输出的图像表示为 $\bm z=h_{\bm\theta}(\bm x)\in\mathbb R^d$，$d$ 是特征维数。这里有一个信息论意义上的硬约束：如果目标类别在这个表示空间里本来就是可分的，那几张标注主要只是用来确定“哪片区域对应哪个类”；可一旦两种不同病害的图像被编码成了同一个向量，信息已经丢失，那么任何只读取这个向量的任务头——无论多复杂——都不可能再把它们区分开。所以评价表示要做控制变量实验：固定来源训练样本、骨干、增强和目标划分，在同一个冻结表示上分别试线性分类、近邻、类别均值这几种简单任务头，差异才能归因于表示；若同时既换表示又换算法，就说不清是谁的功劳。

这里必须插入一个常被元学习论文“比下去”的朴素基线：普通的监督预训练（先在来源上做常规分类训练，再拿这个表示迁移）。它绝不能省，因为大量证据表明，调好的普通预训练加简单任务头，往往不逊于精巧的少样本方法。Chen 等人在统一实现下重新比较了一批少样本分类方法，并专门检查普通微调的跨领域基线；Tian 等人则考察“冻结表示后拟合线性分类器”以及自蒸馏的作用。这些研究给我们的提醒是双向的：在把基准上的差距归功于“元学习”之前，必须先把表示质量和实现细节这些变量控制住；但反过来，也不能因为简单基线很强就断言“所有目标都只需固定特征”——确实有些目标需要调整表示。另外自蒸馏多了一轮教师—学生训练，其机制见第~\ref{sec:ket-integration-student-training}节，这笔准备费用也要记账\scite{ket-chen2019closer,ket-tian2020embedding}

如果准备目标明确设定为“只靠少量支持图像就能分类”，那我们可以反过来塑造表示：专门训练它，让“同类的几个支持点”足以充当查询样本的参照。最简单的参照是什么？用每类支持表示的平均点当“这一类的典型样子”。\term{原型网络}（Prototypical Networks）正是这么做：把每类支持表示的均值当作该类的代表点（原型），再看查询点离哪个原型近，就把类别概率分给谁。对任务$k$的$N$个候选类，令非空$S_{k,c}$为类别$c$的支持子集。采用平方欧氏距离时，代表点与概率写为
\begin{align}
  \bm c_{k,c}
  &=\frac{1}{|S_{k,c}|}
    \sum_{(\bm x_i,y_i)\in S_{k,c}}h_{\bm{\theta}}(\bm x_i),
    \label{eq:ket-transfer-prototype}\\
  p_{\bm{\theta}}(y=c\mid\bm x,S_k)
  &=\frac{\exp\{-\|h_{\bm{\theta}}(\bm x)-\bm c_{k,c}\|_2^2/\tau\}}
    {\sum_{c'=1}^{N}\exp\{-\|h_{\bm{\theta}}(\bm x)-\bm c_{k,c'}\|_2^2/\tau\}}.
    \label{eq:ket-transfer-prototype-probability}
\end{align}
其中$\bm c_{k,c}\in\mathbb R^d$是类别原型，$\tau>0$为温度，标准形式可取1。负距离越大，表示查询越近；指数归一化将$N$个候选的分数转成和为一的概率。选择均值也有具体依据：$\sum_i\|\bm z_i-\bm c\|_2^2$对$\bm c$的导数为$2|S_{k,c}|\bm c-2\sum_i\bm z_i$，置零给出上述均值，且二次项保证这是最小点。它回答的是怎样以单点减少类内平方偏离，不保证多峰类别可由一个点充分代表\scite{ket-snell2017-prototypical}

原型网络最关键的洞见是训练方式：准备阶段就反复“模拟”未来少样本使用的完整场面。做法叫回合式（episodic）训练——每一轮从历史类别里抽出 $N$ 个类、每类抽 $K$ 个支持样本组成一个小任务，再从与支持隔离的剩余图像里抽查询；编码器同时处理这两组，用支持算出原型、让查询按式~\eqref{eq:ket-transfer-prototype-probability}分类，然后只根据查询答得对不对来更新编码器。这样训练，相当于让表示在“被当原型用”的真实场景里被优化，而不是只优化普通分类。具体以非空查询集上的平均负对数概率为损失：
\begin{equation}
  L_k(\bm{\theta})
  =-\frac{1}{|Q_k|}\sum_{(\bm x,y)\in Q_k}
  \log p_{\bm{\theta}}(y\mid\bm x,S_k).
  \label{eq:ket-transfer-episodic-loss}
\end{equation}
梯度同时流过查询表示和支持均值去更新 $\bm\theta$，从而把表示塑造成“适合按支持原型分类”的样子。实现上有个易错点：编码器每更新一步，原型就得用当前的支持表示重新算一遍，不能一直沿用准备开始时算好的旧均值。经过对大量重采历史任务的训练，最终保存下来的是编码器以及“用什么距离、什么温度”这套约定；至于未来新类别的原型，则要等真正遇到该类时，由它自己的支持集现场建立——原型网络并不预存未来类的原型。

用一个均值点代表一整类，隐含假设是“这类大致围成一团”。但同一类花可能有多种形态（全开的、半开的、不同颜色的），取均值会把这些模式平均到一个谁都不在的“空白中心”，反而误导。\term{匹配网络}（Matching Networks）于是放弃单点，改为保留逐个支持样本作为参照，让和查询越相似的支持样本，在投票时标签权重越大。令$f_{\bm\theta}(\bm x,S_k)$与$g_{\bm\theta}(\bm x_i,S_k)$为维度相容的查询和支持表示，$\operatorname{sim}$为规定的相似度，值越大越相似。先把相似度归一化成支持权重，再对同标签权重求和：
\begin{align}
  s_i(\bm x,S_k)
  &=\operatorname{sim}\bigl(f_{\bm{\theta}}(\bm x,S_k),
       g_{\bm{\theta}}(\bm x_i,S_k)\bigr), \nonumber\\
  a_i(\bm x,S_k)
  &=\frac{\exp\{s_i(\bm x,S_k)/\tau\}}
  {\sum_{j=1}^{|S_k|}\exp\{s_j(\bm x,S_k)/\tau\}}, \nonumber\\
  p_{\bm{\theta}}(y=c\mid\bm x,S_k)
  &=\sum_{i=1}^{|S_k|}a_i(\bm x,S_k)\,\mathbf 1\{y_i=c\}.
  \label{eq:ket-transfer-matching}
\end{align}
这里 $i$ 遍历所有支持样本，指示函数 $\mathbf1\{y_i=c\}$ 只把属于类别 $c$ 的那些权重留下相加，温度 $\tau>0$ 控制注意力有多集中（越小越“只盯最像的那一个”）。原论文采用余弦相似度，并进一步让每个样本的编码依赖整个支持集：用双向循环网络读一遍支持序列、让查询通过多次注意力去读取支持表示。这意味着同一张图的编码会随“它和谁同台出现”而变化，表达力更强，但也带来排列处理和循环计算的开销，记账时不能再把它当成一个彼此独立的固定编码器\scite{ket-vinyals2016matching}

将式~\eqref{eq:ket-transfer-matching}代入同一查询损失，即可按历史任务训练匹配网络。固定独立表示时，逐例匹配的参照数量随支持样本数增长，原型则只保留每类一个点，减少存储与比较次数，但改变了类别形状假设。每类一例且编码器、比较分数相同，均值就是该例，两种规则才一致；余弦与平方距离不同，或只有一方用上下文编码时，不能据样本数宣称等价。第~\ref{sec:ket-transfer-references}节进一步计算目标参照，并检查多例情形的差异。

\subsection{可迁移语义联系的准备}
\label{sec:ket-transfer-semantic-preparation}

原型和匹配都需要新类的图像来求参照，可零样本场景里连图像都没有，只有一句话描述，这时连原型都无从算起。破局的思路是引入一层“属性”作为中间桥梁：既然描述（“花瓣细长、叶缘有锯齿”）说的是可观察特征，那就先学会从图像中辨认这些特征，再用特征去匹配描述。这里两类知识各司其职——类别描述回答“应该找什么”，来源监督回答“怎样从像素里把它认出来”；只有把这两者接起来，文字描述才能真正成为新类预测的依据。若没有这座桥，所谓零样本就只是给输出换了个名字而已。

先用 $m$ 个二值属性把这座桥具体化：类别 $c$ 的描述就是一个 $m$ 维 0/1 向量 $\bm a_c\in\{0,1\}^m$，第 $j$ 位表示该属性是否存在。\term{直接属性预测}（direct attribute prediction, DAP）分两步：先在来源数据上训练 $m$ 个属性预测器 $q_j(\bm x)\approx p(a_j=1\mid\bm x)$（即“这张图有第 $j$ 个特征吗”），再拿预测出的属性向量去和目标类别的描述比对。这个设计有个很实际的优点：保存下来的只是 $m$ 个属性预测器，输出维数与目标类别数完全无关，所以新增类别不需要重训网络，只要补一段新描述。来源监督可以是逐张图标注属性，也可以偷懒把整类的描述赋给该类每张图；但后者默认了“同类对象都具备这些属性”，遇到类内变异、遮挡或图像里根本看不到的属性（比如“根有毒”）时就站不住\scite{ket-lampert2009attributes}

属性预测要变成类别分数，还须说明概率模型。在理想化DAP模型中，类别$c$确定$\bm a_c$，即$p(\bm a\mid y=c)=\mathbf1\{\bm a=\bm a_c\}$，并假定给定属性后，类别与图像条件独立。记类别先验为$\pi_c=p(y=c)$，且$p(\bm a_c)>0$。先由贝叶斯公式求属性对应的类别概率，再对不可直接观察的属性求和，得到
\begin{equation}
  \begin{aligned}
    p(y=c\mid\bm a)
      &=\frac{\pi_c\,\mathbf 1\{\bm a=\bm a_c\}}{p(\bm a)},\\
    p(y=c\mid\bm x)
      &=\sum_{\bm a\in\{0,1\}^m}
          p(y=c\mid\bm a)\,p(\bm a\mid\bm x)\\
      &=\pi_c\,\frac{p(\bm a_c\mid\bm x)}{p(\bm a_c)}.
  \end{aligned}
  \label{eq:ket-transfer-attribute-link}
\end{equation}
类别描述固定以后，求和只剩$\bm a=\bm a_c$的一项。到这里，等式描述的仍是理想模型中的精确后验，还没有说明怎样从$m$个单属性预测器计算联合属性概率。

为使这个比值可计算，接下来分别近似图像条件下的属性联合分布与属性先验，记$p_j=p(a_j=1)$：
\[
  \begin{aligned}
    p(\bm a_c\mid\bm x)
      &\approx\prod_{j=1}^{m}
        q_j(\bm x)^{a_{c,j}}[1-q_j(\bm x)]^{1-a_{c,j}},\\
    p(\bm a_c)
      &\approx\prod_{j=1}^{m}
        p_j^{a_{c,j}}(1-p_j)^{1-a_{c,j}}.
  \end{aligned}
\]
两个乘积分别假定条件属性独立与先验属性独立，前者不推出后者：混合不同图像后，属性仍可能相关。将它们代入式~\eqref{eq:ket-transfer-attribute-link}，便能只用$q_j$、$p_j$和目标描述计算类别分数：
\begin{equation}
  s_c(\bm x)=\pi_c
  \prod_{j=1}^{m}
  \frac{q_j(\bm x)^{a_{c,j}}
        [1-q_j(\bm x)]^{1-a_{c,j}}}
       {p_j^{a_{c,j}}(1-p_j)^{1-a_{c,j}}},
  \label{eq:ket-transfer-attribute-score}
\end{equation}
分子比较预测属性与类别描述，分母把这份支持与属性原先的出现频率相比。例如，某项属性的先验概率为$0.1$，图像预测它存在的概率为$0.8$，对应的属性比值就是8。这表示图像对该属性的支持高于原先频率；比值本身不是类别概率。

实际计算还需要估计这些概率。目标属性先验通常未知，采用来源估计后，式~\eqref{eq:ket-transfer-attribute-score}应视为分类分数，不能继续当作理想模型下的精确后验。Lampert等按来源类别属性表逐列取均值，不按各类图像数量加权。本节采用相同逐类口径，再加对称伪计数$\delta>0$避免分母为零：
\[
  \hat p_j=\frac{\sum_{c\in C_S}a_{c,j}}{|C_S|},
  \qquad
  p_j=\frac{\sum_{c\in C_S}a_{c,j}+\delta}{|C_S|+2\delta}.
\]
这里$C_S$非空，$\hat p_j$为未平滑的逐类均值，$p_j$为平滑后的实际计算值，必在$(0,1)$内；图像多的类别不会获得额外权重。对非空目标候选集，本节取均匀$\pi_c=1/|C_T|$，不以最终查询标签估计频率。连乘可改为累加对数，零预测概率则须按预定裁剪或极限约定处理。来源先验、属性独立性和预测器误差都可能偏离目标实际情况，所以将分数归一化也没有自动完成概率校准。

如果来源已经有类别概率模型，还可以先预测源类，再按源类属性表汇总属性。Lampert等将这条路径称为间接属性预测（indirect attribute prediction, IAP）：
\[
  p(a_j=1\mid\bm x)
  =\sum_{c\in C_S}a_{c,j}\,p(y=c\mid\bm x).
\]
这把$q_j$的学习换成对源类概率的复用，仍受源类与属性覆盖限制。两条路径都不能突破描述本身的区分能力：新种甲、乙的属性向量及先验相同，DAP分数就相同；两者仅在叶缘上不同，而照片没有叶片，观测也可能无法区分。增加描述词语只有在它们对应可观察、可学习的差异时才提供新依据。

人工列出的属性词表终究有限，很难描述任意新类；连续的图文表示则学习一种更通用的比较接口。\term{对比语言图像预训练}（Contrastive Language--Image Pre-training, CLIP）同时训练一个图像编码器和一个文本编码器，用对比目标让“真实配对”的图文在共同空间里尽量靠近、不配对的尽量远离。经过海量图文对训练，这个空间天然就携带了视觉—语义联系，类别不再只是一个固定编号，而是可以用任意一句话描述的位置。完整对比训练细节见第~\ref{part:efficient-knowledge-use}部分；迁移时保存下来的是两个编码器、共同空间和尺度约定，目标端只需把新类描述编码成一个可与图像比较的参照向量\scite{radford2021}

不过连续空间也带来新的局限：它不要求每个维度对应一个可解释属性，所以你无法像读属性表那样读出“它靠的是花瓣还是颜色”；相似度也未必能分辨否定（“有斑点”与“无斑点”可能很像）、数量或细粒度形态差别。使用时仍要把类别名称、属性、解释当作需要管理的目标描述或样本，记录提供者、挑选模板用的验证材料、以及预训练覆盖范围。还有一个常被夸大的说法要警惕：想声称模型识别的是“来源未见类”，必须真正在来源里排除过该类；一个在海量网页图文上预训练的模型，只是“没在你的目标训练集上做监督拟合”，绝不等于它从没见过这个概念——它很可能早就见过。

\subsection{参数初始化与任务先验的准备}
\label{sec:ket-transfer-initialization-prior}

少量新类图像当然也能拿来做梯度优化，但样本太少时它们根本约束不动动辄百万、千万的全部参数——硬训只会过拟合。来源侧能在这件事上提供两类帮助：一是给更新一个好的起点（初始化），二是直接规定“不同任务的参数大概长什么样”的分布（任务先验）。这两者层次不同：初始化 $\bm\theta$ 是一组确定的数值，从它出发只有一条轨迹；任务先验 $p(\bm w_k\mid\bm\eta)$ 则描述参数的整个可能取值分布，由共享配置 $\bm\eta$ 控制，它不仅有中心，还有方差、相关性甚至多个峰。把先验均值拿来当初始化是可以的，但这一步只传递了“中心”，并没有把方差、相关性或多峰结构这些不确定性信息交给更新过程。下面先讲起点如何按“适应后的损失”来学（MAML 一系），再讲分布如何被目标证据修正（层次贝叶斯一系）。

\subsubsection{参数初始化的学习}
\label{sec:ket-transfer-maml}

这里有一个反直觉但关键的认识：在“新任务只允许走几步梯度更新”的前提下，源任务上的最优参数未必是好起点。因为我们要的不是“起点本身在各任务上表现好”，而是“从这个点出发、走几步之后表现好”——最优解可能处在一个狭窄谷底，附近的新任务几步走不到好地方；反而是一个略次、但周围地形平缓、朝多个方向都容易下降的点更适合当公共起点。\term{模型无关元学习}（model-agnostic meta-learning, MAML）就直接把“支持集更新之后、在查询集上的损失”当作准备目标，从历史任务中学这样一个共享起点。注意“模型无关”并非毫无限制，它只适用于由可微参数刻画、能通过梯度过程训练的模型，结构、损失与任务协议仍需事先选定\scite{finn2017}

对非空支持、查询集，记$L_{S_k}(\bm w)$和$L_{Q_k}(\bm w)$为平均损失。固定$M$个历史任务批次，先沿支持损失的负梯度做一步任务内更新，再平均各任务更新后的查询损失以调整共享起点：
\begin{align}
  \bm w_k^{(0)}&=\bm{\theta}, \nonumber\\
  \bm w_k^{(1)}&=\bm{\theta}
       -\alpha\nabla_{\bm{\theta}}L_{S_k}(\bm{\theta}),
       \label{eq:ket-transfer-maml-inner}\\
  J(\bm{\theta})&=\frac{1}{M}\sum_{k=1}^{M}
       L_{Q_k}(\bm w_k^{(1)}(\bm{\theta})), \nonumber\\
  \bm{\theta}&\leftarrow\bm{\theta}-\beta\nabla_{\bm{\theta}}J(\bm{\theta}).
       \label{eq:ket-transfer-maml-outer}
\end{align}
其中$\alpha>0$为任务内步长，$\beta>0$为跨任务步长，$\bm w_k^{(1)}$是支持更新后的任务模型参数。各任务从同一个当前$\bm\theta$独立复制，不能让下一任务继承上一任务的末态。外层以降低适应后的平均查询损失为目标，不保证每次更新都下降；保存的是经验证选出的共享起点，而非把各任务分类器依次拼成一个分类器。

训练 MAML 时有个技术核心必须讲清：外层要优化的是起点 $\bm\theta$，但查询损失是在“内层更新之后的参数” $\bm w_k^{(1)}$ 上算的，所以查询损失对起点的依赖要先穿过内层那一步更新，绝不能简单地把查询点处的梯度直接当成起点梯度（那是 FOMAML 的近似，稍后讲）。在损失足够光滑的假设下，令 $\bm H_k=\nabla^2L_{S_k}(\bm\theta)$ 为支持损失在起点处的 Hessian，$\bm I$ 为同维单位阵，对内层更新式用链式法则求导，得到：
\begin{align}
  \frac{\partial\bm w_k^{(1)}}{\partial\bm{\theta}}
     &=\bm I-\alpha\bm H_k, \nonumber\\
  \nabla_{\bm{\theta}}J(\bm{\theta})
     &=\frac{1}{M}\sum_{k=1}^{M}
       (\bm I-\alpha\bm H_k)^{\top}
       \nabla_{\bm w}L_{Q_k}(\bm w_k^{(1)}).
       \label{eq:ket-transfer-maml-chain}
\end{align}
直观上，查询梯度告诉我们“在适应后的参数处该往哪走”，而雅可比 $\bm I-\alpha\bm H_k$ 负责把这个要求“换算”回起点处该怎么调整——它刻画了“起点动一点，经过一步更新后终点会动多少”，这正是元学习目标和直接在查询上训练的根本差别。二阶连续可微时 Hessian 对称；实际实现不必存储参数维数平方那么大的 Hessian 矩阵，只需计算 Hessian 与向量的乘积即可。当然，深度网络里有 ReLU 这类非光滑点，此时求导约定由具体实现决定，不能据此声称网络处处满足上面的经典二阶光滑条件。

多步内层更新会逐步改变起点到任务末态的映射。对预先确定、在此求导中不依赖模型参数的步长$\alpha_t$，令$\bm J_k^{(t)}=\partial\bm w_k^{(t)}/\partial\bm\theta$，则
\[
  \bm J_k^{(0)}=\bm I,\qquad
  \bm J_k^{(t+1)}
  =\bigl[\bm I-\alpha_t\nabla^2L_{S_k}(\bm w_k^{(t)})\bigr]
    \bm J_k^{(t)}.
\]
因此第$k$项的起点梯度为$(\bm J_k^{(T)})^\top\nabla L_{Q_k}(\bm w_k^{(T)})$，各步雅可比按实际更新顺序相乘；反向自动微分沿同一轨迹反传。普通混合任务训练直接降低共同模型参数上的合并训练样本损失，没有先形成任务解再评价的映射。两者即便使用相同骨干与优化器，也应在相同支持证据和更新预算后比较，更新前准确率回答不了起点是否利于适应。

算法~\ref{alg:ket-transfer-maml}将重置、适应和外层反馈串成完整训练过程。模型的任务预测输出接口须能表达历史与目标标签，初始模型参数的生成规则也须预先规定。目标使用只执行冻结方案中的支持更新，历史采样和跨任务求导均属于准备费用。

\begin{algorithm}[MAML的准备与目标适应]
  \label{alg:ket-transfer-maml}
  \begin{algorithmic}[1]
    \Require 历史训练任务、独立元验证任务；内层步数$T$与步长$\alpha_t$；外层步长$\beta$
    \Ensure 冻结的共享初始化及给定目标支持集的适应模型
    \State 初始化共享模型参数$\bm{\theta}$
    \While{训练预算尚未耗尽，且未满足元验证停止条件}
      \State 抽取$M$个历史任务，各自构造不重叠的$S_k,Q_k$
      \For{$k=1,\ldots,M$}
        \State 复制共享起点：$\bm w_k^{(0)}\gets\bm{\theta}$
        \For{$t=0,\ldots,T-1$}
          \State $\bm w_k^{(t+1)}\gets\bm w_k^{(t)}
            -\alpha_t\nabla L_{S_k}(\bm w_k^{(t)})$
        \EndFor
        \State 用$\bm w_k^{(T)}$计算查询损失$L_{Q_k}$
      \EndFor
      \State 沿内层更新轨迹求导，用平均查询梯度更新$\bm{\theta}$
    \EndWhile
    \State 按元验证结果选定并冻结初始化、步长与步数
    \State 新目标从该初始化出发，仅在$S_T$上执行约定的$T$步更新
    \State 对$Q_T$预测输入产生预测；隐藏答案只进入独立评分
  \end{algorithmic}
\end{algorithm}

下面用一个极简一维例子，亲手算一算“好起点”和“适应后更好”为什么会指向不同位置，从而理解 MAML 到底在优化什么。
\begin{example}[同一组分类任务中的起点与适应后损失]
  \label{ex:ket-transfer-maml-numeric}
  为了只比较起点与一步适应，构造一维花卉特征$z$，用$wz$作为二值标签$y\in\{-1,1\}$的预测分数，按$\operatorname{sign}(wz)$分类，并以平方损失训练。这不是MAML实测结果。任务1支持样本为$(1,1),(-1,-1)$，查询为$(2,1),(-2,-1)$；任务2支持为$(2,1),(-2,-1)$，查询为$(1,1),(-1,-1)$。这些数值对代表独立样本，支持与查询幅度刻意不同，以展示准备目标对评价位置的依赖。

  对每组采用$\frac{1}{2n}\sum_i(wz_i-y_i)^2$，得到
  \[
    L_{S_1}(w)=\tfrac12(w-1)^2,\quad
    L_{S_2}(w)=\tfrac12(2w-1)^2,
  \]
  \[
    L_{Q_1}(w)=\tfrac12(2w-1)^2,\quad
    L_{Q_2}(w)=\tfrac12(w-1)^2.
  \]
  对两个支持损失求和并令导数为零，得普通混合训练起点$w_{\mathrm{pre}}=3/5$。固定同一步长$\alpha=1/4$，内层更新分别给出$w_1^{(1)}=3w/4+1/4$和$w_2^{(1)}=1/2$。将这两个末态代入查询损失并取平均，元目标为
  \[
    J(w)=\tfrac14(3w/2-1/2)^2+\tfrac1{16},
  \]
  其最小点是$w_{\mathrm{meta}}=1/3$。在普通起点$w=3/5$处，$\mathrm dJ/\mathrm dw=0.3$。任务2内层雅可比为$1-\alpha\cdot4=0$，说明一步后模型参数不再依赖起点，因此它的起点梯度为零。直接将适应后查询梯度作用于起点，会漏掉这个依赖关系。

  更新前，普通与元起点的平均查询损失依次为$0.05$与$5/36\approx0.1389$，普通起点较小；同样一步后，依次为$0.1025$与$0.0625$，元起点较小。分数符号在此均正确，准确率不能区分它们，平方分数损失却能。这个构造只说明固定更新预算下应评价适应后目标，不是实际花卉任务中MAML优于预训练的证据。
\end{example}

精确回传上面的二阶梯度，需要在内存里保存或反复重算内层更新的计算图，这在大网络上很贵。\term{一阶MAML}（first-order MAML, FOMAML）做了个干脆的近似：把那个雅可比 $\bm I-\alpha\bm H_k$ 直接当成单位阵，也就是忽略“终点随起点的曲率变化”，直接拿适应后参数处的查询梯度去调起点。这样省掉了穿过支持梯度的二阶求导，但要注意它省的只是二阶项：历史任务采样、内层更新、查询反馈这套元训练流程一个都不少，所以它仍是元学习，只是被改变的梯度不再是精确的元梯度。

Reptile 走得更简：它甚至不显式划分查询集、也不计算外层梯度，而是直接利用“在任务内训练若干步后参数相对起点移动的那段位移”。每个历史任务都从同一个共享起点出发，独立地用普通梯度下降走若干步得到末态 $\bm w_k^{(T)}$，然后让共享起点朝“各任务末态的平均位置”移动一小步，步幅由 $\epsilon>0$ 控制：
\begin{equation}
  \bm{\theta}\leftarrow\bm{\theta}
   +\epsilon\frac1M\sum_{k=1}^{M}
        \bigl(\bm w_k^{(T)}-\bm{\theta}\bigr).
  \label{eq:ket-transfer-reptile}
\end{equation}
Reptile不要求另取查询集计算外层梯度，方向来自任务内小批次更新的累计位移。一步普通SGD时，它退化为混合任务梯度方向；多步则在已更新模型参数上计算后续梯度，保留了任务内小批次的相互作用。FOMAML仍读显式查询梯度，故两者不是同一种去掉Hessian的算法。Reptile的独立元验证和最终目标查询隔离仍须保留\scite{ket-nichol2018reptile}

为看清位移中的相互作用，记同一任务两批损失在共同起点的梯度为$\bm g_1,\bm g_2$，第二批Hessian为$\bm H_2$。若第二批梯度可作带二阶余项的局部泰勒展开，固定小步长$\alpha$时有
\[
  \bm w^{(2)}-\bm{\theta}
    =-\alpha(\bm g_1+\bm g_2)
      +\alpha^2\bm H_2\bm g_1+O(\alpha^3).
\]
第二批梯度取在$\bm\theta-\alpha\bm g_1$，展开后产生$\bm H_2\bm g_1$项，而非仅在原起点相加两个梯度。Reptile实现仍只调用一阶梯度；这里用Hessian解释小步位移，没有要求算法额外执行二阶求导。

还可以控制内层范围，检查是否必须改变表示。将模型参数分为骨干$\bm\theta_b$与任务头$\bm\theta_h$，\term{几乎无内循环}（Almost No Inner Loop, ANIL）只在任务内更新头，外层仍训练骨干与头起点。查询损失直接依赖查询表示，也经支持表示依赖头适应，故精确ANIL没有在整个准备期间冻结骨干。原研究在MiniImageNet等指定网络与少样本任务中支持特征复用的解释；若新病害特征已被骨干丢失，目标端只调头仍无法恢复它\scite{ket-raghu2019anil}

MAML 显式地只走固定几步，若想让目标把支持损失训练到接近最优点，就得保存一条很长的更新轨迹，这可能成为准备阶段的主要开销。换个思路：与其对“有限步更新”这个过程求导，不如改一下内层目标，让任务的最优解始终被“拴”在共享中心附近，然后直接对这个最优解满足的驻点条件求导（隐函数求导），这样就不必保存到达解的整条轨迹。\term{隐式MAML}（implicit MAML, iMAML）在内层目标里加入一个以共享参数 $\bm\theta\in\mathbb R^P$ 为中心的二次惩罚项；对固定支持集 $S$ 和 $\lambda>0$，假定最小值可达并选定一个任务解，记为：
\begin{equation}
  \bm w^\star(\bm\theta)
  \in\argmin_{\bm w\in\mathbb R^P}
    \left\{L_S(\bm w)
       +\frac{\lambda}{2}\|\bm w-\bm\theta\|_2^2\right\}.
  \label{eq:ket-transfer-implicit-inner}
\end{equation}
这里$\lambda$控制支持拟合与接近共享中心的取舍。即使内层迭代很多步，惩罚中心仍是$\bm\theta$，与只在起点使用$\bm\theta$不同。式~\eqref{eq:ket-transfer-implicit-inner}也改变了原先固定步数的适应映射，不能把它视为只更换求导实现而保持MAML目标完全相同。

设所用局部任务解唯一、随$\bm\theta$可微，$L_S$二阶连续可微，并令$\bm H_S=\nabla^2L_S(\bm w^\star)$。当$\bm H_S+\lambda\bm I$可逆时，对内层驻点条件求导，有
\begin{align}
  \nabla L_S(\bm w^\star)
       +\lambda(\bm w^\star-\bm\theta)&=\bm0,\nonumber\\
  (\bm H_S+\lambda\bm I)
       \frac{\partial\bm w^\star}{\partial\bm\theta}
       &=\lambda\bm I.
  \label{eq:ket-transfer-implicit-jacobian}
\end{align}
这个导数由解附近的曲率决定，不要求保存到达该解的全部中间模型参数。以$L_S(w)=(w-1)^2/2$、$\lambda=1$为教学例，任务解为$(1+\theta)/2$，对共享中心的导数为$1/2$。它既不等于一阶近似中的1，也不等于无正则支持目标充分优化后对起点的零导数。

外层仍以隔离查询的损失$L_Q(\bm w^\star)$选择共享中心。记$\bm g_Q=\nabla_{\bm w}L_Q(\bm w^\star)$，先解$P$维线性方程，再得到这项查询损失的元梯度：
\begin{equation}
  \begin{aligned}
    (\bm H_S+\lambda\bm I)^{\top}\bm v&=\bm g_Q,\\
    \nabla_{\bm\theta}L_Q(\bm w^\star(\bm\theta))
      &=\lambda\bm v.
  \end{aligned}
  \label{eq:ket-transfer-implicit-gradient}
\end{equation}
系数矩阵正定时，可用共轭梯度和Hessian向量积近似求$\bm v$，无须显式构造完整Hessian；仅有可逆性时还须选择适用的线性求解器。实际任务模型参数只近似满足驻点条件，线性方程也只近似求解，两项误差都进入元梯度。较小的方程残差能否带来较小导数误差，还依赖曲率条件，不能由省存储推出精确性或泛化保证\scite{ket-rajeswaran2019-imaml}

适应解也可以只是一项凸任务头。MetaOptNet将支持表示交给线性支持向量机等凸学习器，查询损失再通过任务头最优条件传回表示模型参数。局部解分支可微、相应Karush--Kuhn--Tucker（KKT）条件的线性化非退化时，可以计算这种隐式导数；活跃约束切换等位置须按具体求解条件处理。这里内层凸不表示神经表示的完整准备目标凸。与ANIL按有限步更新头相比，改变的是任务头的求解及求导对象；目标端仍须付出求解费用，并保留支持与查询的用途边界\scite{ket-lee2019-metaoptnet}

\subsubsection{任务参数先验的学习}
\label{sec:ket-transfer-prior}

同一个少样本任务，能很好拟合这两张病害图的参数其实有无穷多组，一个孤零零的起点完全没有表达“这些合理解通常集中在哪、分散程度多大”。层次先验模型则用一个共享的先验分布去概括“历史上各任务的参数是怎样分布的”，等目标证据一来再用贝叶斯方式修正这个分布——比单点初始化保留了丰富得多的不确定性信息。设 $\bm\eta$ 为共享先验分布配置，$\bm w_k$ 为任务模型参数；本节给定各支持输入，仅对标签建立似然。假定给定 $\bm\eta$ 后任务模型参数独立，给定各 $\bm w_k$ 后任务标签按似然独立，联合分解为
\begin{equation}
  p(\bm\eta,\{\bm w_k,S_k\}_{k=1}^{M})
   =p(\bm\eta)\prod_{k=1}^{M}
      p(\bm w_k\mid\bm\eta)
      \prod_{(\bm x_i,y_i)\in S_k}
      p(y_i\mid\bm x_i,\bm w_k).
  \label{eq:ket-transfer-hierarchical-model}
\end{equation}
式左的 $S_k$ 只是对“支持标签及其已给定输入”的简写，并没有额外去建一个输入本身的生成模型。那么共享先验配置 $\bm\eta$ 怎么从历史任务里学？关键是把未知的任务参数 $\bm w_k$ 积分掉：对每个候选先验，计算它给该任务全部支持标签赋予的边缘概率（即把先验允许的所有参数取值都考虑进来），哪个先验让历史任务的标签在平均意义上最可能出现，就选哪个。形式上，在预先约定的配置候选范围内最大化这些边缘概率的对数和：
\begin{equation}
  \hat{\bm\eta}
  =\argmax_{\bm\eta}\sum_{k=1}^{M}
     \log\int p(S_k\mid\bm w_k)\,
          p(\bm w_k\mid\bm\eta)\,\mathrm d\bm w_k.
  \label{eq:ket-transfer-empirical-bayes}
\end{equation}
积分保留每个先验允许的任务模型参数，而非先为各任务任选一个最优点；对数把独立任务概率的乘积转为求和。用历史训练样本估计先验分布配置称为经验贝叶斯（empirical Bayes），另一个选择是保留$\bm\eta$的后验。两者都需要规定分布族、积分或优化近似与任务关系；一个任意正则系数不能据名称就成为已学到的先验。先验结构仍由独立历史验证任务选择，目标查询答案不参与上述估计。

准备阶段定出 $\hat{\bm\eta}$ 后，到了目标端就是标准的贝叶斯更新：把目标支持集提供的似然与学到的先验相乘，得到参数后验。多数情况下我们只需要后验里的一个代表点，那就取后验最大点（MAP），等价于把“负对数似然”和“负对数先验”两项加起来一起最小化——后者正是一个正则项：
\begin{align}
  p(\bm w_T\mid S_T,\hat{\bm\eta})
  &\propto p(S_T\mid\bm w_T)p(\bm w_T\mid\hat{\bm\eta}), \nonumber\\
  \hat{\bm w}_T
  &=\argmin_{\bm w}\left[
       -\sum_{(\bm x_i,y_i)\in S_T}\log p(y_i\mid\bm x_i,\bm w)
       -\log p(\bm w\mid\hat{\bm\eta})\right].
  \label{eq:ket-transfer-map}
\end{align}
第一行描述完整模型参数分布，第二行给出最大后验（maximum a posteriori, MAP）点；写$\argmin$不保证优化已找到全局最小值。要让模型参数不确定性进入预测，还需计算$\int p(y\mid\bm x,\bm w)p(\bm w\mid S_T,\hat{\bm\eta})\,\mathrm d\bm w$。它将各模型参数的预测输出按后验加权，不能以一个MAP点预测输出代替。

为计算先验强度的作用，暂将目标换成花卉连续测量的简化线性回归，不把类别编号当连续值。设$\bm w_k\in\mathbb R^d$服从$\mathcal N(\bm\mu,\sigma_w^2\bm I)$，$y_i=\bm z_i^\top\bm w_k+\varepsilon_i$，噪声独立服从$\mathcal N(0,\sigma_y^2)$，两方差严格为正。历史任务提供中心$\bm\mu$和差异尺度，噪声尺度亦按预定方案给定或估计。目标$n$个支持特征按行组成$\bm Z\in\mathbb R^{n\times d}$，观测组成$\bm y\in\mathbb R^n$。将似然与先验的二次项配方，后验协方差和均值为
\begin{align}
  \bm\Sigma_T
    &=\left(\sigma_y^{-2}\bm Z^\top\bm Z
                     +\sigma_w^{-2}\bm I\right)^{-1}, \nonumber\\
  \bm m_T
    &=\bm\Sigma_T\left(\sigma_y^{-2}\bm Z^\top\bm y
                     +\sigma_w^{-2}\bm\mu\right).
  \label{eq:ket-transfer-gaussian-posterior}
\end{align}
因$\sigma_w^{-2}\bm I$严格正定，矩阵可逆；$\bm\Sigma_T\in\mathbb R^{d\times d}$，$\bm m_T\in\mathbb R^d$。所得后验为$\mathcal N(\bm m_T,\bm\Sigma_T)$，此时均值恰好也是MAP点。支持样本不仅移动中心，还改变各模型参数方向的方差；单独保存起点$\bm\mu$无法表达这个更新后的分布。

\begin{example}[先验怎样改变少样本估计]
  \label{ex:ket-transfer-prior-numeric}
  构造一维$z_i=1$、噪声方差为1、先验$w_T\sim\mathcal N(0,1)$，支持观测为$2,4$。无先验的似然估计是均值3；按式~\eqref{eq:ket-transfer-gaussian-posterior}，后验精度即方差的倒数为$1+2=3$，均值为$6/3=2$，方差为$1/3$。仅将先验方差减为$1/4$，均值变为$6/(4+2)=1$，方差为$1/6$；将其增为4，均值为$6/(1/4+2)=8/3$，方差为$4/9$。强先验让本模型内的后验更集中，也把中心拉得更近；若真实目标模型参数偏离历史中心，这种集中不表示估计更准确。
\end{example}

Grant等给出MAML的层次贝叶斯解释，将有限步任务估计与跨任务准备联系到经验贝叶斯近似。联系依赖估计与近似方式，不是任意起点都已指定完整模型参数先验\scite{ket-grant2018hierarchical}

这个区别可在一维二次目标中核对。令负对数似然的有效项为$L(w)=h(w-\hat w)^2/2$，$h>0$是曲率，$\hat w$是其最小点。从$\mu$以固定步长$\alpha$做$T\ge1$步梯度下降，递推给出$w^{(T)}=(1-a)\hat w+a\mu$，其中$a=(1-\alpha h)^T$。若改为加入中心在$\mu$的二次先验惩罚，则
\[
  \argmin_w\left\{L(w)+\frac{\lambda}{2}(w-\mu)^2\right\}
    =\frac{h\hat w+\lambda\mu}{h+\lambda}.
\]
在$0<\alpha h<1$下，$0<a<1$，取$\lambda=ha/(1-a)>0$，这个MAP点恰等于早停模型参数，对应高斯先验$\mathcal N(\mu,\lambda^{-1})$。例如$h=1,\alpha=1/2,T=1$给出$\lambda=1$。等价对象仅是指定损失尺度下的模型参数点，先验精度还依赖曲率、步长与步数；非线性网络中的有限步梯度不因此成为全局MAP求解，更没有给出完整后验。多个当前任务怎样共同设计先验结构，见第~\ref{sec:ket-sharing-prior}节；这里讨论历史准备怎样支持未来估计。

\subsection{参数更新规则的准备}
\label{sec:ket-transfer-learned-optimizer}

MAML 学的是“从哪出发”，但出发之后怎么走同样关键：给定一个支持梯度，该迈多大、要不要带动量、不同方向是否用不同步长——这些选择决定了同样几步更新能走出截然不同的轨迹。手工优化器（SGD、Adam）的规则是固定的，那能不能让历史任务来“学习”这套更新规则本身，使它天生就擅长在少样本下快速收敛？这就是学习的优化器（learned optimizer）的思路。设任务模型参数为$\bm w_k^{(t)}$，更新器模型参数为$\bm\psi$，内部状态$\bm s_k^{(t)}$保存此前反馈。先按既定规则创建模型起点与必要的新头，再初始化更新器状态，随后用支持损失的梯度产生增量：
\begin{align}
  \bm s_k^{(0)}
    &=\operatorname{Init}_{\bm\psi}(\bm w_k^{(0)}), \nonumber\\
  \bm g_k^{(t)}
    &=\nabla_{\bm w}L_{S_k}^{(t)}(\bm w_k^{(t)}), \nonumber\\
  (\Delta\bm w_k^{(t)},\bm s_k^{(t+1)})
    &=U_{\bm\psi}(\bm g_k^{(t)},\bm w_k^{(t)},
             L_{S_k}^{(t)},\bm s_k^{(t)}), \nonumber\\
  \bm w_k^{(t+1)}
    &=\bm w_k^{(t)}+\Delta\bm w_k^{(t)}.
  \label{eq:ket-transfer-learned-update}
\end{align}
这里 $L_{S_k}^{(t)}$ 是第 $t$ 步所选支持小批次的损失，$\bm g_k^{(t)}$ 是它对参数的梯度，更新器 $U_{\bm\psi}$ 输出一个与参数同形状的增量 $\Delta\bm w_k^{(t)}$ 以及更新后的内部状态。究竟把损失、梯度、参数、状态中的哪些喂给更新器，由具体实现决定。要特别分清两类在目标端待遇完全不同的量：任务参数 $\bm w_k$ 会被目标反馈持续改变，它承载“当前这个任务的预测”；更新器参数 $\bm\psi$ 则在历史上学好后于目标端冻结，它规定的是“该如何学习”。二者是两项不同的迁移资产，不能混为一谈。

Andrychowicz等以循环网络实现学习更新器。其逐坐标设计让一个小型长短期记忆网络（long short-term memory, LSTM）共享权重，但每个模型参数坐标保留自己的历史状态，并计算该坐标的模型参数增量。这样不必用一个全连接网络读取完整模型参数向量；梯度符号和幅度的预处理也帮助处理不同量级。历史梯度可以经状态影响本步增量，类似可学习的动量与步长，不过训练出这种行为不意味着继承手工优化器的凸问题收敛保证\scite{ket-andrychowicz2016-learning}

状态初始化也是保存接口的一部分。逐坐标$\operatorname{Init}_{\bm\psi}$可以置零或复制学得初态，每个新任务与独立重复都按同一规则重新建立。任务内支持小批次之间沿用状态，任务间不自动传递末态；如果要跨任务保留状态，须另行声明连续适应协议。更新器只能读取许可的支持反馈，最终查询答案不进入状态或模型参数更新。

训练更新器 $\bm\psi$ 时，优化目标不是“让它模仿某个手工步长”，而是直接评价“它在若干任务上走出的整条轨迹好不好”——因为我们真正在乎的是几步之内查询表现。把更新过程展开 $T$ 步，用一组非负权重 $\omega_t$ 指定哪些时刻的查询预测需要做好（比如只要求终点好，或中途也要稳），再对 $M$ 个历史任务取平均：
\begin{equation}
  J(\bm\psi)=\frac1M\sum_{k=1}^{M}
       \sum_{t=1}^{T}\omega_t L_{Q_k}(\bm w_k^{(t)}(\bm\psi)),
       \qquad \omega_t\geq0.
  \label{eq:ket-transfer-optimizer-objective}
\end{equation}
只令$\omega_T$非零，评价的是终点；多个非零权重还约束中途表现，权重须预先规定并至少有一个为正。少样本学习用独立查询评价适应后的泛化，纯函数优化也可直接用所优化的函数值，不能将两者混成同一目标。反向传播沿任务模型参数与更新器状态训练$\bm\psi$；截断轨迹或停止对梯度预测输入再求导，会省计算也会改变得到的梯度，不再是原展开目标的精确导数。

Ravi与Larochelle把少样本分类器模型参数与门控循环状态相联系，学习保留与更新的幅度。核心写成$\bm w^{(t+1)}=\bm f^{(t)}\odot\bm w^{(t)}-\bm i^{(t)}\odot\bm g^{(t)}$，其中$\odot$为逐坐标乘法，$\bm f^{(t)}$为保留门，$\bm i^{(t)}$为更新门，维度均与$\bm w$相同。损失、梯度和循环历史影响这些门，初态还可参与训练，从而联合准备起点与规则。查询标签用于历史外层评价，不交给本任务更新器\scite{ket-ravi2017optimization}

VeLO进一步用更广的神经网络优化任务训练更新器，并利用任务描述与训练进度调整优化行为。论文中无需逐任务搜索学习率的使用方式，仍依赖更新器准备、模型选择和执行成本，不能解释为全过程免调参。换成更大模型、更长轨迹、不同损失或梯度尺度时，要另测外推；短轨迹终点较低，不保证长时间训练稳定\scite{ket-metz2022velo}

费用也分为共享与任务状态。保存更新器需$|\bm\psi|$个模型参数；若任务有$P$个模型参数，每坐标循环状态宽度为$h$，状态另占$O(Ph)$存储，每步还增加更新器前向计算。少做若干次学习率搜索可节省目标端费用，却未必抵消历史训练、共享存储和逐步执行开销，应按实际任务数量核算。

\subsection{任务条件化机制的准备}
\label{sec:ket-transfer-conditioning-preparation}

前面几条路线都默认“要用支持样本就先做梯度更新”。但支持图像其实可以绕开梯度，直接进入预测计算：模型读入这几张“输入—标签”对，在内部形成一个描述“当前任务长什么样”的状态，再带着这个状态去回答查询。\term{任务条件化}（task conditioning）指的就是预测显式地依赖于本次任务给出的证据——共享参数规定了“怎么读证据”，而读出的任务状态随证据即时改变。它和学习更新器的本质区别在于：这个过程可以根本不产生一套优化后的任务参数，没有“训练新模型”这一步，全靠一次前向读取完成；准备阶段要训练的，是“从证据到条件预测”的那个读取映射。

\subsubsection{支持摘要的学习}
\label{sec:ket-transfer-support-summary}

用一个语义固定的回归或分类任务就能说清这种映射。\term{条件神经过程}（Conditional Neural Process, CNP）把每个支持“输入—标签”对先编码成一个向量，再聚合成一个描述当前任务的摘要，最后由这个摘要和查询输入共同产生预测分布。对非空$S_k$，编码器$e_{\bm\phi}$返回$d_r$维表示向量，简单均值聚合为
\begin{align}
  \bm r_k&=\frac1{|S_k|}\sum_{(\bm x_i,y_i)\in S_k}
                   e_{\bm\phi}(\bm x_i,y_i), \nonumber\\
  p_{\bm\phi}(y\mid\bm x,S_k)
      &=p_{\bm\phi}(y\mid\bm x,\bm r_k).
  \label{eq:ket-transfer-cnp}
\end{align}
用简单均值聚合有个明确的取舍：好处是摘要 $\bm r_k\in\mathbb R^{d_r}$ 与支持样本的排列顺序无关（这正是“过程”应有的置换不变性），代价是平均可能抹掉“有几个样本、是否有多个峰”这类细节；若任务确实需要这些，就得换更复杂的聚合方式。解码器方面，回归任务可输出均值和一个正的方差（顺带表达不确定性），分类任务则输出预定标签空间上的概率。要注意标签接口问题：如果未来类别数或标签含义会变，就不能死守一个固定维度的分类输出，必须另行规定标签编码方式。共享参数 $\bm\phi$ 在历史任务上用条件负对数似然训练，到目标端冻结，只用新支持现场算出摘要 $\bm r_T$。

原始CNP训练可从一组观测抽上下文，再预测包含上下文和未观察点的目标集。本章若评价支持之外的泛化，应另报告隔离$Q_k$上的损失，上下文重构不能替代它。编码器和解码器规模固定时，简单均值方案处理$n$个支持与$q$个查询约需$O(n+q)$次点级计算，固定支持摘要可缓存；这项复杂度不自动适用于全点注意力等扩展。

基本因子化解码器在固定支持条件下独立产生查询预测输出，因而没有表示完整的查询相关性。以截距固定、斜率仍不确定的直线为例，对两个正数预测输入，较大的斜率会同时抬高两处的预测值，所以这两处预测输出应共同变化。分别给出每处的均值和方差，只说明各自可能落在哪里；将两处当作条件独立预测输出，没有表达由共同未知斜率带来的联系。

除了同一支持集下的查询相关性，还可以检查不同支持集之间的概率关系：加入一条观测后，新分布是否等于原分布按贝叶斯公式条件化所得的结果。基本CNP的条件分布不保证满足这种关系，因此不能仅凭预测输出均值和方差就认定模型已经给出同一随机函数先验下的贝叶斯预测。方差是否校准是另一项要求，仍需真实查询检验\scite{ket-garnelo2018cnp}

\subsubsection{任务先验下的预测拟合}
\label{sec:ket-transfer-pfn}

上一节讲过，给定任务先验后可以求后验、再做后验预测。但如果每来一个新任务都要重新做一遍积分或近似推断，这笔重复计算开销不小。一个聪明的办法是“摊销”：干脆在准备阶段从指定先验出发大量随机生成学习任务，训练一个单一的条件预测网络，让它学会“看到任意一组支持样本，就直接吐出查询的预测分布”——把无穷多次任务专属推断，压缩成一次前向。\term{先验样本拟合网络}（Prior-Data Fitted Network, PFN）就是这条路线。但务必记住它的贝叶斯身份是有条件的：它拟合的是“你所指定的那个任务先验”所生成的后验预测，并不是只要采用了上下文接口就自动成了贝叶斯方法；先验设错，它学到的“后验预测”也跟着错。

沿用第~\ref{sec:ket-transfer-prior}节的线性回归设定，先抽任务模型参数$\bm w\sim p_0$，再生成支持样本$S$及查询样本$(\bm z,y)$。本段直接把回归特征$\bm z$作为预测输入；其抽样律预先固定且不依赖$\bm w$，标签满足$y=\bm z^{\top}\bm w+\varepsilon$。给定$\bm w$后，支持与查询独立生成，查询标签只用于准备损失。于是$p_0$和样本生成规则共同确定了$p_0(y\mid\bm z,S)$，不需要先收集无穷多个真实植物园。

令$\Phi$为允许的共享模型参数集合，每个$\bm\phi\in\Phi$确定一个条件预测密度$q_{\bm\phi}(y\mid\bm z,S)$。反复抽取上述任务及样本，最小化隐藏查询标签的负对数密度，得到先验样本负对数似然（prior-data negative log-likelihood）：
\begin{equation}
  \min_{\bm\phi\in\Phi}J_{\mathrm{PFN}}(\bm\phi),
  \qquad
  J_{\mathrm{PFN}}(\bm\phi)
    =\mathbb E_{\bm w,S,\bm z,y}
      [-\log q_{\bm\phi}(y\mid\bm z,S)].
  \label{eq:ket-transfer-pfn-objective}
\end{equation}
期望遵循先抽任务、再抽任务内样本的共同生成规则。分类时$q_{\bm\phi}$表示类别概率；这里的连续回归采用密度，密度可以大于1，负对数值也不必非负。对固定$S,\bm z$，生成查询标签的条件律正是$p_0(\cdot\mid\bm z,S)$。因此，在有关期望有限且预测密度覆盖该生成条件密度的支撑时，可将目标改写为
\begin{equation}
  J_{\mathrm{PFN}}(\bm\phi)
   =C+\mathbb E_{S,\bm z}
      \!\left[\operatorname{KL}\!\left(
        p_0(\cdot\mid\bm z,S)
        \,\|\,q_{\bm\phi}(\cdot\mid\bm z,S)
      \right)\right],
  \label{eq:ket-transfer-pfn-kl}
\end{equation}
其中$C$是生成条件律的平均熵，与$\bm\phi$无关。这一步来自把$-\log q$写成$\log(p_0/q)-\log p_0$后取条件期望。只有模型类能够表示该条件律且优化达到相应最小值时，预测才在先验生成的支持与查询取值上几乎处处等于该后验预测；有限网络和有限准备计算通常只给近似\scite{ket-muller2022-pfn}

这项拟合可以用已有高斯数例检查。取一维$z=1$、先验$w\sim\mathcal N(0,1)$、噪声方差1，支持标签为2和4。前面已经算得模型参数后验$\mathcal N(2,1/3)$；新标签还带有方差1的观测噪声，所以理想查询分布为$\mathcal N(2,4/3)$。PFN应从这组支持样本给出接近该分布的预测，而不是只复制后验均值2，或把模型参数的后验方差$1/3$直接当作新标签方差。目标端固定训练后的$\hat{\bm\phi}$，学习算法把支持样本$S_T$变成预测器$\bm z\mapsto q_{\hat{\bm\phi}}(\cdot\mid\bm z,S_T)$，无需为该任务再次更新共享模型参数。一次计算传入多条查询也不自动成为直推式预测；应检查注意力掩码是否允许查询彼此影响，再按实际依赖声明权限。

先验仍决定上述近似希望达到什么。如果真实目标的函数形状、噪声或特征联系不在生成任务的范围内，增加同一先验下的合成任务不会自动补上这份差异。支持样本数超出准备范围、模型容量有限或优化未完成，也可能使近似变差。把重复推断转为共享模型的离线训练属于摊销思路，可参考\BookChapterRef{chap:pgm-parameter-learning}{卷6《模型：概率图模型》的“参数学习”章}；这里额外固定的是整项学习问题的生成先验，而不是只为单条观测计算潜变量后验。

TabPFN将先验生成任务与上下文预测用于表格分类和回归。Hollmann等2025年版本用大量合成表格任务准备模型，目标端读取特征列、支持标签及待预测行。核心模型的一次条件预测可以前向完成，但论文默认分类和回归方案分别采用4种和8种预处理、后处理组合并平均预测，不等于整项服务只执行一次前向。任务生成、离线训练、开发任务选模和默认集成都要计费；该版本的小中型表格实验范围也不是任意表格任务或后续版本的无条件保证\scite{ket-hollmann2025-tabpfn}

\subsubsection{示范读取规则的学习}
\label{sec:ket-transfer-demonstration-reading}

若任务的输入输出天生就是序列形式（典型如语言），还有最直接的一招：把支持样本当成示范，原样拼进输入序列里。MetaICL 在准备阶段从历史任务抽出“支持输入—标签对”和独立查询，把示范和查询拼成一个序列，用查询正确答案的负对数概率去训练“读懂示范并照着做”的能力。用可选说明扩展这一接口，准备目标写为
\begin{equation}
  \min_{\bm\phi}
  \mathbb E_k\left[
    -\log p_{\bm\phi}
       (y\mid D_k,(\bm x_1,y_1),\ldots,(\bm x_n,y_n),\bm x)
  \right],
  \label{eq:ket-transfer-metaicl}
\end{equation}
式中期望包括历史任务及其支持、查询采样，最小化在预定模型参数范围内进行。末尾$(\bm x,y)$为查询，$y$只进入损失，不作为预测自身的预测输入；序列答案的对数概率由词元条件概率求和。$D_k$是本书为“说明加示范”保留的可选扩展，原始MetaICL核心是示范与查询，并不要求说明。准备后，未见任务通过读取新示例改变预测而不更新权重，这是显式训练上下文适应的路径\scite{ket-min2022metaicl}

也可以不专门训练，直接把一个现成的大语言模型当条件预测器：GPT-3 式的少样本使用就是把示例放进提示、而不做针对目标的梯度微调。但要分清它的“准备”是什么——主要是大规模语言建模，而不是上面那种显式的跨任务元训练。由此得到一条重要判据：输入里出现了示例，只能证明当前这次调用使用了示例，既不能据此断定训练阶段用了元学习，也不能把效果全归功于示范里的正确标签（模型可能主要受益于格式或预训练知识）。稳妥的做法是记录来源模型、任务覆盖和目标证据格式，并用标签打乱等控制实验去检验证据到底起了多大作用\scite{brown2020}

无论准备的是表示、语义编码器、先验、更新器还是条件模型，都须保存使用所需版本与初始化规则。它们的训练与验证进入式~\eqref{eq:ket-reuse-amortization}的准备费用$C_{\mathrm{prep}}$。要判断多次使用是否抵消这项投入，可将来源选择并入本章$C_{\mathrm{adapt},k}$，以$C_{\mathrm{predict},k}$记单次预测费用，对$K$个实际目标任务累加：
\begin{equation}
  C_{\mathrm{total}}
    =C_{\mathrm{prep}}
     +\sum_{k=1}^{K}
       \bigl(C_{\mathrm{adapt},k}+q_kC_{\mathrm{predict},k}\bigr),
  \label{eq:ket-transfer-total-cost}
\end{equation}
这里$K$仅表示本段实际任务数，$q_k$是任务$k$的请求次数；若逐次费用不同，应按真实请求逐项累加，而非强用一个常数。检索与状态管理依发生阶段归入适应或预测，所有计算费用须用相同单位，人工和存储另列或按明确单价折算。第~\ref{sec:ket-reuse-resource-benefits}节要求计入失败尝试、匹配无迁移对照，并区分生命周期总投入与利用已有模型的增量决策。较低历史准备损失还不是未来任务收益的证据。

\section{目标任务的适应}
\label{sec:ket-transfer-adaptation}

准备工作就绪后，视角转到目标端：手头已经有一个交付好的来源模型，面对新种甲、新种乙的那点证据，到底怎么用？归纳起来有三条路径，区别在于“被直接改变的对象”是什么。第一条是构建显式参照：把类别描述或支持表示变成可直接比对的参照点，模型权重纹丝不动，纯靠前向比对作答。第二条是组织任务上下文：把说明和示范塞进上下文，让条件化模型在激活和注意力层面被“临时设置”，同样不改权重。第三条才是用目标监督去优化，把新信息持久地写进任务参数。这三者不是必须顺次走完的阶段，可以自由组合——比如先用文本参照给出初步预测，再用少量标注拟合一个校正头；也要澄清，前两种方式虽然不更新网络权重，但它们显然仍然在使用目标的类别描述或样本，并不是“没用目标信息”。

\subsection{预测参照的构建}
\label{sec:ket-transfer-references}

先看只给新种甲、乙描述的情形。最自然的约定是把预测范围限定为“只能在这两个指定的未见类里选”，这就是传统零样本学习（zero-shot learning, ZSL）的候选条件。但真实部署时，镜头里既可能出现新种，也可能冒出月季、鸢尾这些已见类，模型不该被迫把一切都归到新种里；一旦已见类也允许出现在测试中，就进入广义零样本学习（generalized zero-shot learning, GZSL）。这时若 $C_T$ 专指未见类，候选集合要扩成 $C_S\cup C_T$，而且必须分别报告两部分错误（把已知误判成新类、把新类误判成已知）。要特别小心：GZSL 里已见类享有充分训练的天然优势，候选数量和 softmax 归一化也都变了，所以在纯 ZSL 上调好的阈值不能直接照搬。

描述足以区分类别时，DAP直接使用式~\eqref{eq:ket-transfer-attribute-score}；CLIP类模型则将文本参照和查询图像编码为非零向量，再分别归一化成单位向量$\bm t_c,\bm z\in\mathbb R^d$。为在当前$N$个候选类中比较，采用点积相似度并以温度$\tau>0$归一化：
\begin{equation}
  p(y=c\mid\bm x,D_T)
  =\frac{\exp(\bm z^\top\bm t_c/\tau)}
         {\sum_{c'=1}^{N}\exp(\bm z^\top\bm t_{c'}/\tau)}.
  \label{eq:ket-transfer-text-reference}
\end{equation}
单位向量的点积就是余弦相似度，越大表示在学来的共同空间里越接近。要分清这里两样东西的来源：“参照写了什么”来自目标描述，而“怎样算相似、怎样把相似变成概率”的机制来自来源侧的图文训练。输出的只是候选集合上的模型概率，并不会自动成为经过校准的后验概率。还有个容易忽略的归一化效应：候选集里加入一个新类，会改变 softmax 的分母，从而改变所有原有类的概率——哪怕编码器一个参数都没动。用多个提示模板时可以按预定规则把它们的编码平均、再重新归一化，但挑选模板所用的验证材料和多次编码的费用都要记账。

\begin{example}[描述、支持样本与原型的比较计算]
  \label{ex:ket-transfer-reference-numeric}
  构造两类描述与向量，分别计算语义和图像支持的作用，数值不代表实际植物属性或模型实测结果。教学属性依次为“花瓣细长”“叶缘有锯齿”，甲类描述$(1,0)$，乙类$(0,1)$。查询属性预测$(0.8,0.3)$，两类先验均为$1/2$；来源属性表各列有一半类别取1，逐类均值与对称平滑均给出$1/2$。DAP公共因子消去后，甲、乙分数正比于$0.56$与$0.06$，分别归一化为$28/31\approx0.9032$、$3/31\approx0.0968$。

  另一组共同空间中，单位文本参照为$\bm t_{\text{甲}}=(1,0)$、$\bm t_{\text{乙}}=(0,1)$，单位图像表示为$\bm z=(0.8,0.6)$。取$\tau=0.2$，两类相似度经缩放为4和3，甲类概率为$e^4/(e^4+e^3)\approx0.7311$。两次都选甲，概率却来自不同假设与表示，不能比较数值大小来判断哪种方法更可靠。

  再让目标提供每类两张标注图像，以另一个冻结视觉编码器得到甲类$(0,0),(0,2)$、乙类$(2,0),(2,2)$，查询为$(0.5,1)$。原型为$(0,1)$、$(2,1)$，查询平方距离为$0.25,2.25$，温度1给甲概率$1/(1+e^{-2})\approx0.8808$。同表示上逐例匹配的距离为$1.25,1.25,3.25,3.25$，相同负平方分数求和也得到这个概率。这次相等来自对称点的共同偏移，不是多样本时两规则恒等，更不与前两种证据条件等价。
\end{example}

要检验均值丢掉了什么，换成一维教学支持：甲为$-2,2$，乙为$0.5,0.7$，查询为2。原型0与0.6的平方距离是4与1.96，故单原型选乙；温度1的逐例匹配给甲总权重$1+e^{-16}$，乙总权重$e^{-2.25}+e^{-1.69}$，故选甲。甲的两端模式平均到了没有支持点的位置，说明减少参照改变了判别规则。保留逐例或增加类内原型可表达更多模式，但极少支持也可能不足以估计这些模式。

前面的比较都默认“答案一定在候选类里”，可目标一旦含私有类，候选内比较就缺少一个“退出机制”：当图像谁都不像时，模型需要被允许不选任何已知类。部分重叠任务要先识别出共享部分 $C_{\cap}$，开放集任务则要允许未知样本不被强行匹配到已知参照。最简单的退出判据是用最大类别概率当“接受分”：越接近 1 越像某个已知类，只有当它超过预定阈值 $\gamma$ 才给出类别，否则判为未知，即
\begin{equation}
  \hat y(\bm x)=
  \begin{cases}
    \argmax_{c\in C_{\mathrm{known}}}p(c\mid\bm x),
       & \max_{c\in C_{\mathrm{known}}}p(c\mid\bm x)\geq\gamma,\\
    \mathrm{unknown}, & \text{否则}.
  \end{cases}
  \label{eq:ket-transfer-rejection}
\end{equation}
这里 $\gamma\in[0,1]$，出现并列最大时按预先固定的规则处理。阈值必须由独立验证集、可接受的已知类误拒率或成本约束来定，绝不能拿最终的未知测试答案去把它“选到最佳”（那是泄漏）。而且要明白这个简单规则为什么只是权宜之计：softmax 天生要把概率归一化到总和为 1，于是哪怕没有一个已知类真的合适，它也可能硬塞给某个类一个不低甚至很高的概率，所以“最大概率高”并不总意味着“真的认识”。距离阈值、留出未知验证类、或专门训练一个拒识模型能提供别的依据，但每一样都额外带来描述、样本和假设；仅靠已知类验证，最多能约束误拒率，无法保证对任意未知类都检得出。

分类参照只回答候选类别问题，不能继承为任意预测输出任务。检测可在候选区域上比较参照，但还须产生区域、恢复位置并去除重复；分割须保留空间分辨率，生成须接入序列解码条件。另一方面，冻结特征上拟合线性头虽然也由支持建立分类规则，却用优化形成任务模型参数，属于第~\ref{sec:ket-transfer-parameter-adaptation}节。直接均值与模型参数拟合的区别是如何形成预测器，不是有没有使用标签。

\subsection{任务上下文的组织}
\label{sec:ket-transfer-context}

同一张照片，既可以被问“这是什么物种”，也可以被问“这株病情如何”——所以条件模型必须先从当前证据里弄明白“这次到底要回答哪个问题”。\term{上下文学习}（in-context learning, ICL）指的就是模型利用当前输入中给出的示例，临时形成本次任务的行为。本节强调它的状态性质：共享权重保持冻结，变化只发生在激活、注意力读取或当前任务状态里，是一种“用完即散”的临时设置，并不会自动变成下次会话仍然保留的参数变化。第~\ref{sec:ket-transfer-conditioning-preparation}节讲的是怎么在准备阶段训练出这种读取接口，这里则聚焦目标端该如何组织可被读取的证据。

例如指令规定“按观测样本给出病害复查等级”，教学示范给出“无斑点且边缘完整$\rightarrow$低”“有扩散斑点$\rightarrow$高”，最后放入无答案的新样本。说明定义责任，示范展示对应，查询是待预测对象。把指令改成物种识别，原示范就没有提供新目标的映射。这里使用的是文本观测样本；要检查真实图片还需图像编码接口，文本模型不能直接读取未编码图片。序列边界必须明确，查询答案只保留在评价端。

当当前证据可能同时符合好几条规则时，与其押注最像的一条，不如显式地给每条规则算一个权重、让它们投票。\term{任务后验推断}（task posterior inference）就是这个做法：为潜在规则（或任务变量）$u$ 设定先验和“在该规则下看到这些证据的似然”，再依据支持和说明得到规则的后验，最后对所有规则的预测按后验加权。假设$u$给出标签机制，且查询预测输入不额外改变规则后验，则
\begin{equation}
  p(y\mid\bm x,S_T,D_T)
  =\int p(y\mid\bm x,u)\,
          p(u\mid S_T,D_T)\,\mathrm du.
  \label{eq:ket-transfer-task-posterior}
\end{equation}
积分覆盖所设规则或规则模型参数范围，不是只选一条最高权重规则；离散$u$改用求和。例如候选中可包含“扩散斑点出现就判高等级”。若任务身份还决定预测输入分布，当前查询$\bm x$也可提供任务判别依据，后验须改为$p(u\mid\bm x,S_T,D_T)$。允许联合使用其他查询时，还应按直推协议把它们纳入条件，不能不改条件就扩大可访问样本范围。

不显式建后验时，模型也可以从外部支持样本选择示范或计算读取向量，这称为记忆读取；直接把描述与示范交给网络，则是上下文条件化。任务后验强调概率含义，记忆读取强调证据访问，条件化强调预测依赖，三者可同时存在。例如网络状态可以编码近似后验，注意力也可读取外部支持样本。核算成本与撤销证据时，仍要列清实际保存的是后验概率表或分布配置、样本还是激活，以及缓存多久。

Xie等在文档由混合隐马尔可夫模型生成的设定中，分析从语言建模到文档级潜在概念推断的联系。隐马尔可夫模型让不可直接观察的状态影响序列生成；相应论证还限制模型拟合、提示与预训练分布差异及概念可区分性。上述条件使后验解释可分析，却没有证明任意实际语言模型都学到正确先验或执行精确积分\scite{xie2022bayesianicl}

ICL 究竟“像不像在做学习”？有一种可严格核对的解释：在特定设定下，注意力的一次前向读取，产生的预测竟与“对参数做一步梯度更新”后的预测完全相同。为把对象限定到能算清的范围，取 $d$ 维特征上的线性回归 $\hat y=\bm w^\top\bm z$，对 $n>0$ 个支持样本用平方损失 $L(\bm w)=\frac1{2n}\sum_i(\bm w^\top\bm z_i-y_i)^2$。先求损失梯度、从初始参数走一步步长为 $\alpha$ 的更新，再用更新结果与查询 $\bm z$ 相乘，得到：
\begin{equation}
  \hat y(\bm z)
  =\bm w^{(0)\top}\bm z
   -\frac{\alpha}{n}\sum_{i=1}^{n}
     (\bm w^{(0)\top}\bm z_i-y_i)\,\bm z_i^\top\bm z.
  \label{eq:ket-transfer-context-gd}
\end{equation}
其中支持残差 $r_i=\bm w^{(0)\top}\bm z_i-y_i$ 度量当前模型在第 $i$ 个支持样本上的偏差，决定该样本应贡献多大、什么方向的修正。关键的对照来了：令查询向量 $\bm q=\bm z$、键 $\bm k_i=\bm z_i$、值 $v_i=-\alpha r_i/n$，那么一次“未归一化的注意力读取” $\sum_i(\bm q^\top\bm k_i)v_i$ 恰好就是上式里的那个修正项。注意这里没有 softmax、也不除以权重和，所以权重可正可负，和标准概率注意力并不是一回事。于是，一次前向注意力就能复现一步梯度更新后的同一预测，却不必把结果持久写回参数 $\bm w^{(1)}$。von Oswald 等人在特定编码假设下给出了这种线性注意力构造并研究了简单回归上的训练；但这条代数对应只在受限条件下成立，不能直接外推成“任意 Transformer 都在内部执行梯度下降”的一般结论\scite{ket-vonoswald2023gradient}

这些构造解释某些机制，却不能判断当前模型到底利用了什么。Min等在所研究分类和多项选择任务中，用标签替换等对照观察到标签空间、预测输入分布与格式的作用，不能据此断言任意新映射都不需要正确标签\scite{ket-min2022demonstrations}表~\ref{tab:ket-transfer-context-controls}给出对应检查：固定查询和基本预算，每次只改变一个因素，观察同一目标指标怎样变化。

\begin{table*}[tbp]
  \centering
  \small
  \begin{tabularx}{\linewidth}{p{30mm}XX}
    \toprule
    受控变化 & 保持不变的对象 & 希望区分的解释 \\
    \midrule
    将示范标签随机替换 & 支持预测输入、标签词表、序列长度 & 是否依赖正确预测输入与标签对应，还是主要利用格式与标签范围 \\
    将类别名改为新符号并给出一致映射 & 预测输入和正确对应关系 & 是否学会当前映射，还是沿用预训练中的类别名称习惯 \\
    交换示范顺序 & 全部预测输入与标签对与查询 & 顺序或邻近位置是否主导结果 \\
    仅改变格式和分隔符 & 任务语义与所有示范内容 & 是否对表示接口敏感，而不是任务证据不足 \\
    更换任务指令或移除示范 & 查询预测输入与允许的观测或标注来源 & 任务说明和示范各自贡献什么，二者冲突时依赖谁 \\
    \bottomrule
  \end{tabularx}
  \caption{检查上下文依赖的单因素对照。每行独立执行，固定查询并记录改变因素之外的证据与长度；改名须在支持和正确评价映射中一致，错误标签实验则有意破坏对应。对照本身不预设结果，提示选择只使用验证任务，最终查询答案不参与选择。}
  \label{tab:ket-transfer-context-controls}
\end{table*}

目标端检索可以按查询相似性选择示范，但库中只能有许可样本，不能取回查询自身或按评价单位应隔离的同源答案。较长序列增加编码与生成费用；因果语言模型能缓存固定前缀的键值，缓存却占存储，示范变化还须重算相关部分。跨请求记忆要明确任务切换与清空规则，避免测试任务相互污染。这里仅处理被动证据；强化学习的任务推断还需主动选择动作获得反馈，并把探索代价与后续回报一起评价，可参考\BookChapterRef{chap:rl-transfer}{卷7《学习范式：强化学习》的“任务泛化与适应”章}。

\subsection{模型参数的调整}
\label{sec:ket-transfer-parameter-adaptation}

第三条路径让支持标签通过优化，形成会被持久保存的任务参数。讨论方法前必须先把两个常被搅在一起的选择拆开，否则方法名会误导人。第一个维度是“允许更新多少来源已有参数”：完全不动、只动一部分、还是全部动。第二个维度是“要不要为任务另建一批新参数”：不新增、或新增。两个维度一交叉，很多方法的位置就清楚了：只训练一个新分类头＝不更新已有参数但新增；全量微调再加新头＝全部更新且新增；LoRA 外加解冻归一化量＝部分更新且新增。下面按“全量调整、只调已有部分、附加任务参数”三种机制展开，它们是讲解重点而非互斥的类别，实际方案常是组合。

\begin{table}[htbp]
  \centering
  \small
  \begin{tabularx}{\linewidth}{p{28mm}XX}
    \toprule
    已有模型参数更新范围 & 不新增任务模型参数 & 新增任务模型参数 \\
    \midrule
    不更新 & 直接使用或上下文条件化已有模型 & 新建头、适配器、LoRA、软提示 \\
    部分更新 & 调整已有头、偏置、归一化量或部分层 & 新建头并解冻高层；LoRA并调整已有归一化量 \\
    全部更新 & 预测输出语义兼容时全量微调原模型 & 新任务头与全部来源模型参数共同微调 \\
    \bottomrule
  \end{tabularx}
  \caption{分别记录已有模型参数的更新范围与任务模型参数的新增。每格是一种模型参数来源组合，表中方法名仅代表常用配置，具体冻结与更新仍须按实现说明。新增模型参数少不保证来源全部冻结，也不保证反向传播路径短。}
  \label{tab:ket-transfer-parameter-axes}
\end{table}

模型参数范围也不决定来源怎样准备。普通预训练、MAML、先验估计和学习更新器，可以在接口相容时与不同范围结合。使用学习更新器，须按式~\eqref{eq:ket-transfer-learned-update}重建当前任务状态，并让规则只作用于选定模型参数。比较从相同支持证据出发，记录初始化、更新集合、步长与停止条件；配置依许可的独立验证确定，不能用最终标签选择。

\subsubsection{全量参数的调整}
\label{sec:ket-transfer-full-finetuning}

如果新物种要求连来源表示本身都改变（比如必须学会关注一种来源里没强调过的细节），就需要\term{全量微调}（full fine-tuning）：放开全部来源参数，让它们随目标监督一起调整。输出语义兼容时，目标模型直接从来源参数出发 $\bm\theta_T^{(0)}=\bm\theta_S$；若要新类别或新任务头，则要先规定头的结构和初始化方式，再让它和来源骨干一起训练。要澄清一个措辞：这种情况下来源骨干是被全量调整的，而新头是从零另学的，不能把那个来源里根本不存在的新头也说成“继承来的参数”。

固定单标签新物种分类，用$p_{\bm\theta_T}(y\mid\bm x)$表示目标类别概率。要提高非空支持集正确类别的概率，同时限制允许的模型参数变化，可在预定任务模型范围内最小化
\begin{equation}
  L_T(\bm{\theta}_T)
   =-\frac1{|S_T|}\sum_{(\bm x_i,y_i)\in S_T}
       \log p_{\bm{\theta}_T}(y_i\mid\bm x_i)
       +\lambda\Omega(\bm{\theta}_T),
  \label{eq:ket-transfer-finetuning-loss}
\end{equation}
第一项是平均交叉熵，$\Omega$是预先规定的正则项，$\lambda\ge0$控制其权重；若取自先验，应交代与负对数先验及损失尺度的对应。按训练小批次执行$T$步优化，依据许可验证保存检查点。没有目标验证标签，就使用历史验证任务预先确定的配置，不能临时以查询集选步数。改为检测还须更换位置与匹配损失，单标签交叉熵本身不评价对象坐标。

要判断来源起点有无帮助，应固定目标结构、支持、增强与测试，比较来源初始化和随机初始化在相同预算下的适应曲线。充分训练后的最终风险也可比较，但它回答的不是少量更新预算问题。全量调整允许改变表示，也可能拟合少样本中偶然的背景关系；是否出现这种过拟合，取决于独立查询表现，不能仅由模型参数数量判定。

来源有$P$个模型参数时，全量训练通常为它们计算梯度，Adam等优化器还维护模型参数对应的状态，并处理中间激活。逐任务保存完整副本，权重存储约为$P$个模型参数，另加新头及必要状态。步数少不代表显存或总费用低于参照构建。目标风险下降还只说明当前能力，原物种是否仍可识别须另测。

新类别同时伴随拍摄或观测变化时，可结合第~\ref{sec:ket-update-weighting}节的校正及式~\eqref{eq:ket-update-crossmodal}的配对监督，但须重新核对标签联系、目标覆盖和访问权限。所需无标签批次、配对观测或旧训练样本并非每个协议都提供。保留旧能力的责任还需加入第~\ref{sec:ket-update-retention-validation}节的约束与评价，一次目标适应不能代替长期维护检查。

\subsubsection{已有部分参数的调整}
\label{sec:ket-transfer-existing-parameters}

如果目标仍在识别来源就有的月季、鸢尾，只是拍摄背景变了，往往不必放开全部参数——大动干戈既容易过拟合又费算力，可以先只调整原有分类头或上面少数几层。把来源参数切成冻结的 $\bm\theta_F$ 和可训练的 $\bm\theta_A$ 两部分，冻结部分保持来源值不动，只在允许的子集上最小化前面定义的目标损失：
\[
  \min_{\bm{\theta}_A}
  L_T(\bm{\theta}_F=\bm{\theta}_{S,F},\bm{\theta}_A).
\]
候选范围由选定的层或坐标确定，更新只作用于$\bm\theta_A$，其余模型参数保持来源值。原有头、最高层、已有偏置和归一化量都是可选子集。“已有”按来源中是否存在判断，而非按训练占比或部署是否另存副本判断。

不同开放范围对应不同的“可改变能力”，要按需选择：只调头，改的只是固定表示空间里的那条分类边界；连最高层一起放开，还能改变高层特征如何组合；全量调整则允许连早期、底层的特征都变。这里藏着一个表示能力的硬上限：如果后来换成依赖叶片细斑的病害任务，而冻结的最终表示在编码时就已经把这种细微差异抹掉了，那么只在最终表示上加头是无论如何补不回来的——信息在表示里已经不存在。所以开放范围应结合“目标需要什么特征、支持证据能约束多大能力”来定，不能用一句“解冻的层越多越好”代替在同一任务上的实测比较。

冻结仅规定权重不变，并不删除计算。若只训练顶部头，骨干与预处理均固定，可缓存支持特征供重复拟合，但新查询仍需骨干前向。可训练部件在中间时，损失梯度还须经过后续冻结层回传；不求这些层的权重梯度，仍可能需求表示梯度与激活。因此可训练模型参数比例下降，不自动给出同比例显存或推理节省。

ANIL在目标端只适应任务头，却使用跨任务准备过的骨干，说明部分更新与元学习可以组合。如果目标类别数或病害语义要求重新创建$\bm W_T,\bm b_T$，这些新模型参数就不再属于已有头更新。应列出复制、替换、初始化与冻结的操作；新头同时解冻骨干末层，是附加模型参数与已有部分更新的组合。

\subsubsection{附加任务参数的学习}
\label{sec:ket-transfer-added-parameters}

当原类别集合里根本没有新物种时，光调旧头没有对应输出位，最直接的做法是为任务新增一个分类头。固定来源的 $d$ 维编码器 $h_{\bm\theta_S}$，新建 $N$ 个类的权重 $\bm W_T\in\mathbb R^{N\times d}$ 和偏置 $\bm b_T\in\mathbb R^N$，把线性分数经 softmax 归一化成概率：
\begin{equation}
  \bm p(\bm x)
   =\operatorname{softmax}(\bm W_T h_{\bm{\theta}_S}(\bm x)+\bm b_T).
  \label{eq:ket-transfer-new-head}
\end{equation}
第$c$项对应当前类别$c$的概率，目标训练的新增量共$N(d+1)$个。记$\bm z=h_{\bm\theta_S}(\bm x)$，真实标签独热向量为$\bm e_y$；单样本交叉熵对权重与偏置的梯度是$(\bm p-\bm e_y)\bm z^\top$与$\bm p-\bm e_y$，因而支持标签直接改变新头。平方距离原型也可改写成受约束线性分数，却由支持均值形成，不是这样独立优化权重。

只加头有个天花板：它只能在既有最终表示上划边界，若有用的特征在中间层就缺失或需要微调，加头无能为力；可全量微调又太贵、还容易动到通用能力。于是出现参数高效微调（PEFT）这条折中路线：底座权重冻结，只在网络里插入少量可训练的任务模块。\term{适配器}（adapter）让隐藏表示先经过一个很小的瓶颈变换，再把这份增量以残差方式加回原表示。典型的瓶颈结构对 $\bm h\in\mathbb R^d$ 计算：
\begin{equation}
  \bm h'=\bm h+
     \bm W_{\mathrm{up}}\,
       \sigma(\bm W_{\mathrm{down}}\bm h+\bm b_{\mathrm{down}})
       +\bm b_{\mathrm{up}},
  \label{eq:ket-transfer-adapter}
\end{equation}
降维矩阵$\bm W_{\mathrm{down}}\in\mathbb R^{r\times d}$、升维矩阵$\bm W_{\mathrm{up}}\in\mathbb R^{d\times r}$，偏置分别为$r$维和$d$维，$\sigma$是规定的非线性激活，通常$r\ll d$。每个模块新增$2dr+r+d$个模型参数，小增量初始化让起始计算接近底座。目标损失训练新增量，部署保留模块及底座；每词元位置两次乘法约需$2dr$次乘加，另加激活、偏置与残差，多层和长序列会累积开销。含非线性的模块通常不能吸收为原线性矩阵的固定增量。

Houlsby等的配置在每层注意力表示投影后与前馈子层后各插入串行适配器，处于对应底座残差相加和归一化之前；适配器内部残差与底座残差应分别识别。原配置还训练归一化和任务预测输出部分，所以不能只数瓶颈矩阵。复制任务专有归一化量时须计入新增状态，直接调整来源已有归一化量则是部分已有模型参数更新\scite{ket-houlsby2019-adapters}

适配器是非线性的小模块；若目标只是调整某个线性层的权重，还可以用更直接的低秩参数化：与其存一个和原矩阵同样大的完整增量，不如把增量写成两个瘦矩阵的乘积，用很小的秩 $r$ 去近似。\term{低秩适应}（Low-Rank Adaptation, LoRA）固定原权重 $\bm W_0\in\mathbb R^{d_{\mathrm{out}}\times d_{\mathrm{in}}}$，新增 $\bm A\in\mathbb R^{r\times d_{\mathrm{in}}}$ 和 $\bm B\in\mathbb R^{d_{\mathrm{out}}\times r}$（$r$ 远小于两端维度）。对输入 $\bm h\in\mathbb R^{d_{\mathrm{in}}}$，原支路和低秩支路相加为：
\begin{equation}
  \bm h'=\bm W_0\bm h+s\bm B\bm A\bm h,
  \label{eq:ket-transfer-lora}
\end{equation}
其中$s$是固定缩放，原形式常取$s=\alpha_{\mathrm L}/r$，$\alpha_{\mathrm L}$与本章梯度步长不同。新增$r(d_{\mathrm{in}}+d_{\mathrm{out}})$个模型参数，增量秩不超过$r$；这约束的是本次权重变化，不是证明任务语义具有低秩结构，也不使底座$\bm W_0$成为低秩。与非线性适配器相比，LoRA的并行增量保持为线性映射\scite{ket-hu2021lora}

训练仍需能从初态产生梯度。LoRA常随机初始化$\bm A$并将$\bm B$置零，使增量起始为零。记$\bm\delta=\partial L/\partial\bm h'\in\mathbb R^{d_{\mathrm{out}}}$，按乘法链式法则，单样本梯度为
\[
  \frac{\partial L}{\partial\bm B}
       =s\bm\delta(\bm A\bm h)^\top,\qquad
  \frac{\partial L}{\partial\bm A}
       =s\bm B^\top\bm\delta\,\bm h^\top.
\]
首步$\bm B=0$使$\bm A$梯度为零，但只要相应层表示与损失梯度非退化，$\bm B$可以开始更新。若两因子都置零，两个梯度均为零，普通梯度过程无法启动这项增量。初始函数不变与模型参数可学习，必须同时检查。

线性增量允许在任务固定、格式兼容时预先形成$\bm W=\bm W_0+s\bm B\bm A$，之后只运行同形状矩阵。合并约需$O(d_{\mathrm{out}}d_{\mathrm{in}}r)$运算，并非免费。动态任务服务也可保留共享底座及各任务因子，未合并分支每次预测增加$O(r(d_{\mathrm{in}}+d_{\mathrm{out}}))$计算。量化、并发与切换频率影响部署选择；“可合并”只消除指定形态下的额外预测分支，不消除加载和切换费用。

附加参数还可以放在最前端的词元嵌入处，而不是隐藏层或权重矩阵里。\term{软提示}（soft prompt）是一组由目标损失学出来的连续向量，把它们拼到原词元嵌入前面，相当于给模型一段“可学习的开场白”，引导冻结模型偏向当前任务。原词元嵌入 $\bm E\in\mathbb R^{\ell\times d}$ 有 $\ell$ 个位置，新增 $m$ 个提示向量 $\bm P_T\in\mathbb R^{m\times d}$，按行拼成 $[\bm P_T;\bm E]$ 喂给冻结模型，只用目标答案的负对数似然训练这 $md$ 个参数。要强调软提示通常是稠密的连续向量，并不对应任何可读的词，所以它和人工写的文本提示是两回事。Lester 等的提示调优（prompt tuning）主要就是在嵌入端加向量，不应与那种在每层都插入键值前缀的方法用同一口径数参数\scite{ket-lester2021prompt}

软提示也要用监督与反向传播形成。底座冻结时，求软提示向量的梯度仍需穿过模型层；部署保存提示与底座，长度从$\ell$增为$\ell+m$。稠密自注意力的分数计算从$O(\ell^2d)$增为$O((\ell+m)^2d)$，多出的主要项为$O((2\ell m+m^2)d)$，还会占用上下文与缓存。模型参数少不免逐次费用。文本示范直接作为上下文预测输入，没有这一步连续模型参数训练，证据形式与持久状态也不同。

图~\ref{fig:ket-transfer-added-locations}将这四种机制放在前向路径上比较。应观察新增量读入哪种对象：头读取最终表示，适配器变换中间表示，LoRA并行修改线性映射，软提示改变预测输入条件。这个位置说明可调整的计算不同，不能仅据位置给出同任务风险排序。

\begin{figure*}[htbp]
  \centering
  % Four forward-computation mechanisms; all dimensions are in millimetres.
% No external libraries beyond arrows.meta, positioning and calc are needed.
\begin{tikzpicture}[
  x=1mm,y=1mm,
  font=\figurefont\footnotesize,
  text=BookInk,
  >={Latex[length=2mm,width=1.4mm]},
  frozen/.style={draw=BookMuted,fill=BookPaper,line width=0.8pt,
    text width=30mm,minimum height=11mm,align=center,inner sep=2mm},
  train/.style={draw=BookTeal,fill=FigureModel!10,line width=0.8pt,
    text width=30mm,minimum height=11mm,align=center,inner sep=2mm},
  data/.style={draw=BookBlue,fill=FigureInput!10,line width=0.8pt,
    text width=21mm,minimum height=11mm,align=center,inner sep=2mm},
  add/.style={circle,draw=BookInk,fill=white,line width=0.8pt,
    minimum size=6mm,inner sep=0pt},
  flow/.style={draw=BookInk,line width=0.9pt,-{Latex},rounded corners=1mm},
  wire/.style={draw=BookInk,line width=0.9pt,rounded corners=1mm},
  heading/.style={anchor=west,font=\figurefont\small}
]
  \node[anchor=west,text=BookMuted] at (0,24)
    {共同条件：来源已有模型参数冻结；彩色模块：新增并训练的任务模型参数};
  \node[heading] at (0,13) {(a) 新建预测输出头：读取固定表示};
  \node[data] (head-input) at (13,0) {目标预测输入$\bm x$};
  \node[frozen] (backbone) at (61,0) {冻结底座\\$h_{\bm{\theta}_S}$};
  \node[train] (head) at (113,0) {训练新建头\\$\bm W_T,\bm b_T$};
  \node (head-output) at (157,0) {$\hat y$};
  \draw[flow] (head-input.east) -- (backbone.west);
  \draw[flow] (backbone.east) -- (head.west);
  \draw[flow] (head.east) -- (head-output.west);

  \node[heading] at (0,-22) {(b) 适配器：隐藏表示上的瓶颈残差};
  \node (adapter-input) at (10,-43) {$\bm h$};
  \coordinate (adapter-split) at (24,-43);
  \node[train] (down) at (52,-43) {训练降维\\再经非线性$\sigma$};
  \node[train] (up) at (103,-43) {训练升维\\$\bm W_{\mathrm{up}},\bm b_{\mathrm{up}}$};
  \node[add] (adapter-add) at (140,-43) {$+$};
  \node (adapter-output) at (158,-43) {$\bm h'$};
  \draw[wire] (adapter-input.east) -- (adapter-split);
  \draw[flow] (adapter-split) -- (down.west);
  \draw[flow] (down.east) -- (up.west);
  \draw[flow] (up.east) -- (adapter-add.west);
  \draw[flow] (adapter-split) -- (24,-31) -- (140,-31) -- (adapter-add.north);
  \draw[flow] (adapter-add.east) -- (adapter-output.west);

  \node[heading] at (0,-66) {(c) LoRA：线性映射上的并行低秩增量};
  \node (lora-input) at (10,-87) {$\bm h$};
  \coordinate (lora-split) at (24,-87);
  \node[frozen] (base-weight) at (78,-79) {冻结$\bm W_0$};
  \node[train] (low-rank) at (78,-99) {训练$s\bm B\bm A$};
  \node[add] (lora-add) at (140,-87) {$+$};
  \node (lora-output) at (158,-87) {$\bm h'$};
  \draw[wire] (lora-input.east) -- (lora-split);
  \draw[flow] (lora-split) |- (base-weight.west);
  \draw[flow] (lora-split) |- (low-rank.west);
  \draw[flow] (base-weight.east) -| (lora-add.north);
  \draw[flow] (low-rank.east) -| (lora-add.south);
  \draw[flow] (lora-add.east) -- (lora-output.west);

  \node[heading] at (0,-119) {(d) 软提示：预测输入序列中的连续任务条件};
  \node[train,text width=22mm] (prompt) at (15,-136) {训练提示\\$\bm P_T$};
  \node[data,text width=22mm] (tokens) at (15,-155) {词元嵌入\\$\bm E$};
  \node[data,text width=27mm] (concat) at (66,-145.5) {按序列位置拼接\\$[\bm P_T;\bm E]$};
  \node[frozen] (prompt-model) at (116,-145.5) {冻结模型\\词元长度$\ell+m$};
  \node (prompt-output) at (160,-145.5) {$\hat y$};
  \draw[flow] (prompt.east) -- (38,-136) |- (concat.west);
  \draw[flow] (tokens.east) -- (38,-155) |- (concat.west);
  \draw[flow] (concat.east) -- (prompt-model.west);
  \draw[flow] (prompt-model.east) -- (prompt-output.west);
\end{tikzpicture}

  \caption{来源权重冻结时，任务新增量进入计算的四个位置。实线为前向流，训练模块保存新增模型参数，冻结模块沿用来源权重，加号表示同维相加。适配器只画内部瓶颈残差，LoRA两支路共用同一层表示；软提示与词元嵌入按位置拼接，连续提示不等同于现成文本示范。冻结模块仍参与前向，必要时也传递表示梯度。}
  \label{fig:ket-transfer-added-locations}
\end{figure*}
\FloatBarrier

\begin{example}[固定底座下的模型参数与计算核算]
  \label{ex:ket-transfer-parameter-counts}
  固定一个教学底座：12个Transformer层、隐藏宽度$d=768$、新任务5类。下面只核算指定配置的模型参数与注意力大小，不假设各方法风险相同。新线性头有$5(768+1)=3845$个模型参数。

  每层插两个$r=16$的适配器，则每个$2\cdot768\cdot16+16+768=25360$，24个共608640，加新头612485。若每层两个已有层归一化（LayerNorm）的缩放与偏置也可训练，再加$12\cdot2\cdot2\cdot768=36864$，总数649349。LayerNorm按每个位置的隐藏特征归一化，其缩放、偏置是原有量；因此后一配置同时包含已有部分更新。

  每层仅在查询和值两个$768\times768$映射上加$r=8$的LoRA，每矩阵12288个模型参数，共294912，加新头298757。每个未合并分支、每位置约增加12288次乘加；合并后免这条额外分支，但还需核算合并与切换。

  20个软提示增加15360个模型参数，复用现成生成预测输出时无需同一个分类头；若另加上述五类头，合计19205。长度128增至148时，单注意力头分数矩阵从16384增为21904个元素，多5520个，约33.69\%。这个比例只对应分数矩阵，不是完整前向时间或显存增幅。
\end{example}

若未来只允许训练小模块，还可在历史任务上准备它的适应接口。Meta-Adapters保留冻结底座$\bm\theta_S$、任务内可调普通适配器$\bm\phi$，并在其前后加入元适配器$\bm\omega$。历史内层改普通适配器，查询外层训练元适配器、普通适配器起点与逐层步长。为了让小模块经一步支持更新后降低查询损失，目标为
\begin{align}
  \bm\phi_k'
    &=\bm\phi_0-\bm\alpha\odot
       \nabla_{\bm\phi}L_{S_k}
         (\bm{\theta}_S,\bm\phi_0,\bm\omega), \nonumber\\
  \min_{\bm\omega,\bm\phi_0,\bm\alpha}\quad
  &\frac1M\sum_{k=1}^{M}
    L_{Q_k}(\bm{\theta}_S,\bm\phi_k',\bm\omega).
  \label{eq:ket-transfer-meta-adapters}
\end{align}
式中$\bm\alpha$把每层步长广播到该层适配器坐标，外层在规定的模块模型参数与步长范围内优化。目标端固定底座和元适配器，从$\bm\phi_0$出发只调整普通适配器；元适配器仍执行并参与表示梯度传播。因此迁移的既有小模块起点，也有帮助其适应的共享计算。共享$\bm\omega$与逐任务普通适配器须分开存储计数，实际配置的任务头也不能漏算\scite{ket-bansal2022metaadapters}

前面分述的各种方法，到底各自要存什么、适应时算多少、预测时又算多少？表~\ref{tab:ket-transfer-operation-cost} 回到同一个新物种任务，把“保存的任务状态、一次性适应费用、逐次预测费用”三笔账并排列出，便于横向比较。为读懂表中的量级记号，约定：$n$ 为支持样本数，$N$ 为类数，$d$ 为表示维数，$T$ 为更新次数，$b$ 为每次小批量大小，$F$ 为单样本过一遍底座的前向费用。$B_{\mathcal A}$表示单样本反向费用加按样本分摊的模型参数更新费用，$\mathcal A$为可训练模型参数集合，因而$Tb(F+B_{\mathcal A})$只是一种主要项估算；批量硬件效率须另测。序列费用写成$F(\ell)$以显式保留长度，学习更新器与检索索引还须单列。

\begin{table*}[htbp]
  \centering
  \small
  \setlength{\tabcolsep}{4pt}
  \renewcommand{\arraystretch}{1.2}
  \begin{tabularx}{\linewidth}{@{}>{\raggedright\arraybackslash}p{24mm}
      >{\raggedright\arraybackslash}p{36mm}
      >{\raggedright\arraybackslash}X
      >{\raggedright\arraybackslash}X@{}}
    \toprule
    操作 & 保存的任务状态 & 一次性适应费用 & 逐次预测费用 \\
    \midrule
    语义参照 & 属性表或约$Nd$个向量参照分量 & 名称、属性或描述的整理与编码；不必优化模型 & 图像编码与$N$类比较；向量形式约$F+O(Nd)$ \\
    \addlinespace
    支持样本或原型 & 逐例约$nd$个表示分量，原型约$Nd$个；另存标签 & 提取$n$个带标签样本的特征约$nF$，原型汇总另需$O(nd)$ & 逐例约$F+O(nd)$，原型约$F+O(Nd)$ \\
    \addlinespace
    上下文条件化 & 许可的说明、示范与历史，或其摘要、键值缓存 & 检索与组织证据，计算摘要或固定前缀 & 依接口为$F(\ell)$；长上下文与生成另计 \\
    \addlinespace
    预测头拟合 & 已有头或新建头，约$N(d+1)$个模型参数 & 可缓存固定特征时约$nF+O(TbNd)$ & 底座前向与头部计算，约$F+O(Nd)$ \\
    \addlinespace
    全量微调 & 完整任务权重或等价差量，含必要的新头 & 约$Tb(F+B_{\mathcal A})$；$\mathcal A$含全部来源模型参数及新头 & 适应后完整模型前向 \\
    \addlinespace
    附加任务模型参数 & 适配器、低秩增量或软提示模型参数，以及底座版本 & 前向、相关反向与新模型参数更新；LoRA合并和切换另计 & 底座前向加模块或预测输入位置开销；LoRA合并后无额外分支 \\
    \bottomrule
  \end{tabularx}
  \caption{允许使用相同证据时，各操作的状态与主要计算项。方案可组合多行，也不必用尽证据；新头既是头拟合也是附加模型参数。支持比较采用固定独立编码，完整上下文编码需另加费用；头拟合假定可缓存特征，全量项采用单样本分摊估算，LoRA合并要求固定任务与兼容格式。来源编码器、语义模型和元训练另按式~\eqref{eq:ket-reuse-amortization}计入准备，不给出未经实测的性能排序。}
  \label{tab:ket-transfer-operation-cost}
\end{table*}

模型参数量、标签量、训练时间和预测费用回答四个不同问题。软提示模型参数少却须长路径反传，原型不做梯度更新却使用支持标签，全量训练的一次适应费用还可能被大量请求摊薄。判断整体效率，应固定允许使用的证据与目标标准，并把检索、提示组织、长度、更新、延迟与状态存储同列，不能以其中一项替代全部投入。

最后要承认迁移的边界：当目标明显超出历史覆盖时，没有哪种技巧能保证成功，这时更要允许自己换方案。比如来源只学过整图物种、目标却依赖极微小的叶斑，正确顺序是先核对观测和表示里是否还保留这种细节（信息若已丢失，换什么训练技巧都没用），再结合“增加标注、改变更新次数”后的独立风险曲线去定位真正的瓶颈。而且“没有改善”绝不唯一指向表示不足——它同样可能来自监督给错、优化没调好、或语义根本没对齐；逐项排查后再决定是补观测、扩大训练范围，还是干脆回退到从零开始的普通目标训练。判断迁移收益要沿多个预算点、配着匹配的无迁移对照一起比，不能只挑对自己有利的某一步结果报告。等这套能力形成、转入长期服务后，版本、记忆和持续费用的维护，就交给第~\ref{sec:ket-update-maintenance}节。

\section{本章小结}
\label{sec:ket-transfer-summary}

本章可以收束为一条主线：回到新植物园，来源模型或样本之所以有用，前提是它和目标之间确实存在可被利用的联系——而这种联系对不同任务长在不同地方。识别新物种，需要的是类别含义的对齐和把新类区分开的证据；改判病害，需要观测和表示里保留了病斑这类细节；定位花朵，则需要对象级的位置信息和框与对象的匹配结构，这些都不是“特征通用”一句话能覆盖的。还要再次给“未知拒识”降温：它只是在现有类别目录不够用时提供一个体面的退出，绝不会自动完成给新类命名、发现新类内部结构、或建立新分类能力这些后续工作。因此动手迁移前，目标语义、允许使用的观测与标注、以及评价范围必须先界定清楚——只有这样，你才能说清迁移究竟减少了哪一种错误。

准备什么和怎么使用必须一一对应：一个可比较的表示让支持样本能形成参照，图文之间的联系让文字描述得以进入预测，好的起点、先验和更新器帮助参数在少样本下被估准，而条件模型则干脆直接读取任务证据、不产生任务参数。元学习的共性是把“接受任务证据之后表现如何”纳入历史准备，它并不规定目标端一定要走几步梯度——零梯度的上下文适应同样可以是元学习的产物。落到目标端，还要坚持分别记录两个维度：来源已有参数更新了多大范围、又新增了多少任务参数。适配器、LoRA、软提示插入的位置、梯度回传的路径、以及部署时能否合并都不同，所以“某方法参数少”绝不能被当成一条统一的效率结论——参数少的软提示照样要走长路径反传、要付更长序列的注意力开销。

在理论细节上，要记住“适应的梯度长什么样”取决于你采用的学习算法，不能一概而论：MAML 这类有限步更新是对整条计算轨迹求导（含二阶项，FOMAML 才把它近似掉），而 iMAML 这类隐式方法是对“带二次惩罚条件的适应最优解”求导，二者优化的目标、保存的中间量和引入的近似误差都不同，不能混为一谈。PFN 则代表另一种摊销哲学：把“指定任务先验下的后验预测”直接拟合进一个共享条件模型，从而省掉目标端一遍又一遍的推断；但省下重复求解的同时，它仍背着先验设错（错配）、离线准备成本高、以及实际方案常要对多种预处理预测做集成这几笔代价。

验证环节有两道隔离缺一不可：任务内部要隔离支持与查询，任务之间更要隔离历史准备、模型选择和最终测试。要证明“来源确实帮助了学习”，正确做法是固定可用证据、目标标准和预算，在完全独立的目标任务上画出适应曲线，并与从零训练的无迁移对照比较；只报告准备损失低、或某一次结果好看都不足为据。成本要按真实使用次数摊销，而且失败的尝试、反复的选择都得计入账目，不能只算成功那条路径。最后要清楚这份证据的边界：它说明的是“这一次能力如何形成、花了多少成本”，至于旧能力有没有被破坏、未来能否继续学新东西，仍是第10、11章讨论的两个独立问题，需要各自的长期验证来回答。
